\documentclass[11pt]{article}
\usepackage[margin=1in]{geometry}
\usepackage{amsmath,amssymb,amsthm,mathtools}
\usepackage{booktabs}
\usepackage[colorlinks=true,linkcolor=blue,citecolor=blue]{hyperref}

\newtheorem{theorem}{Theorem}[section]
\newtheorem{proposition}[theorem]{Proposition}
\newtheorem{lemma}[theorem]{Lemma}
\theoremstyle{definition}
\newtheorem{definition}[theorem]{Definition}
\newtheorem{remark}[theorem]{Remark}
\newtheorem{finding}[theorem]{Finding}
\newtheorem{openproblem}[theorem]{Open Problem}
\newtheorem{corollary}[theorem]{Corollary}

\newcommand{\ph}{\varphi}
\newcommand{\Q}{\mathbb{Q}}
\newcommand{\R}{\mathbb{R}}
\newcommand{\Z}{\mathbb{Z}}
\newcommand{\Mah}{\mathrm{M}}
\newcommand{\Var}{\operatorname{Var}}
\newcommand{\E}{\operatorname{\mathbb{E}}}
\newcommand{\Cov}{\operatorname{Cov}}
\newcommand{\forcedtag}{{\scshape [forced]}}
\newcommand{\declaredtag}{{\scshape [declared]}}
\newcommand{\computedtag}{{\scshape [computed]}}
\newcommand{\opentag}{{\scshape [open]}}

\title{\bfseries The Field Surprisal Geometry --- Full Worked Edition\\[2pt]
\large Mahler--Gibbs Structure, Dual Flatness, Thermodynamics, Distances,\\
and a Totally Geodesic Two--Statistic Surface}
\author{Echo--Squirrel Research \quad(\texttt{AceTheDactyl})}
\date{2026}

\begin{document}
\maketitle

\begin{abstract}
We attach a canonical positive--definite Fisher--Rao (surprisal) geometry to a number field's
emission catalog and develop it to four further, fully forced layers. The catalog's Mahler
measures $\{2,3,5,\ph,\ph,\ph^4,\ph^4\!-\!1\}$ sum to $Z=17+4\sqrt5$ and, through the description
cost $\log\Mah$, define the \emph{Mahler--Gibbs family} $p_i(\beta)\propto\Mah_i^{-\beta}$, a forced
exponential family. On it we prove: surprisal is affine in cost with slope $\beta$ and Fisher equals
the cost variance (\S4); the geometry is \emph{dually flat}, with the Gibbs curve an $e$--geodesic
and $\mathrm{KL}$ between temperatures a Bregman divergence of $\log Z$ (\S5); its thermodynamics
obey the identity \textbf{$C(\beta)=\beta^2 I(\beta)$}---Fisher information is a heat capacity---with
a distinguished peak--fluctuation temperature (\S6); its Fisher--Rao distances are Bhattacharyya
angles, the seven seeds forming a regular spherical simplex at mutual distance $\pi$, with an
explicit information length (\S7). Adjoining the degree as a second forced statistic yields a
two--parameter field geometry whose Fisher matrix is $\nabla^2 A$; we prove its Gaussian curvature is
\emph{exactly} the constant $\tfrac14$---an isometric round $\tfrac14$--sphere, embedded as a
\emph{ruled} surface (swept by great circles)---because a single--outcome indicator makes the family a
spherical suspension whose metric is a round sphere for any curve (\S10). No such surface is totally
geodesic: $\log\Mah$ alone spans a $6$--dimensional algebra, so none lies in a great subsphere. Generic
and multi--set--indicator statistics are not even constant curvature (controls give $4/21$, $857/16928$,
and sign--changing values), and among the catalog's forced invariants only degree---the single--outcome
indicator---gives constant $\tfrac14$ (\S11), joined by a second over $\ph^4$ from the conjugate--angle
charge; in fact the constant--curvature--$\tfrac14$ surfaces are \emph{exactly eight}---one per seed
plus the golden--pair merge---and no non--indicator statistic gives another (\S12). Two independent
catalogs form a Riemannian product with block--diagonal Fisher and zero cross--curvature (\S13).
Every load--bearing identity is machine--verified
(\texttt{field\_surprisal\_v2.py}, \texttt{suspension\_theorem.py}, \texttt{field\_surprisal\_tier2.py};
$12/12$, $9/9$, $15/15$; exact arithmetic). This edition closes four fronts. The eight--surface classification is proved \emph{in full} by a
finite branch calculus---an exact window identity plus a plane lemma (\S\ref{sec:classdone})---so the
non--indicator exclusion is now \forcedtag{} with no sampling residue. The temperature is dissolved:
product functoriality plus measurable Cauchy \emph{forces} $p\propto\Mah^{-\beta}$, making $\beta$ a
coordinate on a forced $e$--geodesic of finite total information length $L_{\mathrm{tot}}\approx5.6462$
running from the golden merge to the $\ph^4$ vertex (\S\ref{sec:tempcoord}). Iterated indicator joins
give round $\tfrac14$--spheres in every dimension, realized by the catalog triple (cost, degree,
charge parity) (\S\ref{sec:join}); and the first totally geodesic family appears at exactly $k=5$,
uniquely the golden--merged simplex $\R[\log\Mah]$. A fifth front closes in this revision: the
window identity is proved \emph{conceptually}---$P=Z^2\sum_{|s|=4}q_4(s)\,w_s$ by Sylvester's
determinant identity plus Cauchy--Binet (Theorem~\ref{thm:master})---resolving Open Problem (3) and
retiring the branch calculus from the critical path. This revision then closes three of the remaining open fronts: the master identity lifts to every
$k$ (Sylvester $+$ Cauchy--Binet over $(k{+}2)$--subsets), yielding windowed Gauss obstructions, a
rank--$\le1$ window criterion for constant $\tfrac14$, and a complete $k=3$ double--indicator
census---exactly $26$ windowwise--flat classes, the disjoint within--level pairs; the eight--surface
census is proved \emph{catalog--invariant}, with count $\sum_\ell(2^{m_\ell}-1)$ over cost levels
(generic catalogs need no branch analysis at all, and a new branch--rank calculus closes necessity
for four catalogs including the full one); and the temperature selection resolves as a
\emph{dichotomy}---no metric--canonical operating point exists, and the candidate principles
provably disagree, so the declared election is the answer's final form. This revision closes the last two structural fronts. Necessity holds at \emph{every} $k$: constant $\tfrac14$ is exactly windowwise rank $\le1$, because window collisions are squarefree---the Salem square cannot act on windows, only the golden swap can, and twin symmetry neutralizes it---so the $26$--class census is the complete constant--$\tfrac14$ double--indicator landscape and the analytic problem reduces to finite linear algebra. And the compositum \emph{couples}: on its own tensor catalog the Mahler--Gibbs cost is provably non--separable (interaction contrast exactly $\log 2$; golden sector $-6\log\ph$; interaction tensor of exact rank $5$), resolving the multi--field question on its forced horn. This final revision completes the program: the landscape is classified at every $k$---constant $\tfrac14$ is exactly the \emph{partitioned--affine} condition, by a moment trichotomy on window circuits, with catalog counts $8/56/95/31/1$ and the indicator--join conjecture refuted---and the coupled compositum family is proved \emph{not} constant--$\tfrac14$: the coupling is curvature--visible, already at the uniform point where the cross--covariance vanishes exactly, and $\Delta$ is not determined by the charge data. Sixteen harnesses ($61/61$,
$17/17$, $4/4$, $16/16$, $12/12$, $24/24$, $44/44$, $16/16$, $13/13$, $9/9$, $72/72$, $18/18$, $28/28$, $13/13$, $21/21$, $9/9$; exact
arithmetic over $\Q(\log2,\log3,\log5,\log\ph)$, plus $\log7$ for the extended catalog) verify
every new decision.
\end{abstract}

\section{Introduction and discipline}
The companion $\lambda=2c$ manuscript charges a number field's growth rule a description cost
$\log\Mah(\theta)$ and reads its parameter geometry as the (flat) trace form. Here we build the
complementary object: the positive--definite Fisher--Rao geometry of the catalog's outcome
distribution, and four forced layers on top of it. Nothing new is assumed about the field; we use
only its forced data. Each claim is tagged \forcedtag{} (theorem/exact computation), \computedtag{}
(validated numeric), \declaredtag{} (a principled invariance--free choice), or \opentag. All
symbolic identities are checked in exact \texttt{sympy}; the one numeric object (\S10 curvature)
uses a routine validated against four metrics of known curvature.

\section{The catalog and its forced Mahler data}
\begin{table}[h]\centering
\begin{tabular}{lllc}
\toprule
seed & minimal polynomial & $\Mah$ (exact) & degree\\ \midrule
$\sqrt2$ & $x^2-2$ & $2$ & $2$\\
$\sqrt3$ & $x^2-3$ & $3$ & $2$\\
$\sqrt5$ & $x^2-5$ & $5$ & $2$\\
$\ph$ & $x^2-x-1$ & $\ph=(1+\sqrt5)/2$ & $2$\\
$\tau$ & $x^2+x-1$ & $\ph$ & $2$\\
$\ph^4$ & $x^2-7x+1$ & $\ph^4=(7+3\sqrt5)/2$ & $2$\\
$K$ & $x^4+5x^2-5$ & $\ph^4-1=(5+3\sqrt5)/2$ & $4$\\
\bottomrule
\end{tabular}
\caption{The catalog $\Omega$, with $\Mah(\alpha)=\prod_{|\zeta|>1}|\zeta|$ and degrees. All exact
and machine--verified.}
\end{table}

Each measure is $\prod_{|\zeta|>1}|\zeta|$ over the roots $\zeta$. For a quadratic $x^2-D$ ($D>1$)
both roots $\pm\sqrt D$ lie outside the unit circle, giving $\Mah=D$; the golden seeds keep only the
dominant root, $\Mah(\ph)=\Mah(\tau)=\ph$, and $\Mah(\ph^4)=\ph^4$. The quartic $K$ is the one
non--obvious case: $x^4+5x^2-5=0$ gives $x^2=\tfrac{-5\pm3\sqrt5}{2}$, hence real roots $\pm K$ with
$|K|=5^{1/4}/\ph<1$ (inside) and imaginary roots $\pm i\beta$ with $\beta^2=\tfrac{5+3\sqrt5}{2}>1$
(outside); only the imaginary pair contributes, so $\Mah(K)=\beta^2=\ph^4-1$.

\begin{lemma}[Field constant]\label{lem:Z}
$Z:=\sum_{\alpha\in\Omega}\Mah(\alpha)=10+2\ph+2\ph^4-1=17+4\sqrt5.$ \forcedtag
\end{lemma}
\begin{proof}
$2+3+5=10$; $\ph+\ph=1+\sqrt5$; $\ph^4+(\ph^4-1)=2\cdot\tfrac{7+3\sqrt5}{2}-1=6+3\sqrt5$; summing
gives $10+(1+\sqrt5)+(6+3\sqrt5)=17+4\sqrt5$.
\end{proof}
The integer $17$ is a computed coincidence of this catalog, not a structural invariant, and is read
as nothing more.

\section{The Mahler--Gibbs family}
The forced cost is $E_i:=\log\Mah_i$; the maximum--entropy law at inverse temperature $\beta$ is
\[
p_i(\beta)=\frac{\Mah_i^{-\beta}}{Z(\beta)},\qquad Z(\beta)=\sum_j\Mah_j^{-\beta},\qquad
A(\beta):=\log Z(\beta).
\]
This is an exponential family with sufficient statistic $\log\Mah$ (structure \forcedtag; the value
of $\beta$ is \declaredtag). Distinguished temperatures: $\beta=0$ (uniform); $\beta=1$ (MDL/Occam
prior $p\propto1/\Mah$); $\beta=\lambda=\sqrt5$ (the field's forced exchange rate); $\beta=-1$
(exact anchor $p\propto\Mah$, algebraic in $\Q(\sqrt5)$).

\section{Surprisal geometry}
\begin{proposition}[Surprisal--cost affinity]\label{prop:affine}
$S_i(\beta):=-\log p_i(\beta)=\beta\log\Mah_i+\log Z(\beta)$; hence
$\partial S_i/\partial(\log\Mah_i)=\beta$ and $S_i-S_j=\beta\log(\Mah_i/\Mah_j)$. \forcedtag
\end{proposition}
\begin{proof}
From $p_i(\beta)=\Mah_i^{-\beta}/Z(\beta)$,
$S_i=-\log p_i=-\log\Mah_i^{-\beta}+\log Z(\beta)=\beta\log\Mah_i+\log Z(\beta)$. The partial
derivative and the difference are immediate. The slope of surprisal against the Mahler cost is
exactly the inverse temperature $\beta$; at $\beta=\lambda=\sqrt5$ the surprisal therefore equals the
growth--rule cost $\lambda\log\Mah_i$ up to the common constant $\log Z$.
\end{proof}

\begin{proposition}[Fisher $=$ cost variance]\label{prop:fisher}
$I(\beta)=A''(\beta)=\Var_\beta(\log\Mah)\ge0$, where $A(\beta):=\log Z(\beta)$. \forcedtag
\end{proposition}
\begin{proof}
Since $Z'(\beta)=\sum_j\Mah_j^{-\beta}(-\log\Mah_j)$,
\[
A'(\beta)=\frac{Z'(\beta)}{Z(\beta)}=-\sum_j p_j(\beta)\log\Mah_j=-\E_\beta[\log\Mah].
\]
Differentiating once more,
\[
A''(\beta)=-\frac{d}{d\beta}\,\E_\beta[\log\Mah]
=\E_\beta[(\log\Mah)^2]-\big(\E_\beta[\log\Mah]\big)^2=\Var_\beta(\log\Mah),
\]
the standard identity that the second derivative of the log--partition of an exponential family is
the variance of its sufficient statistic; non--negativity is that of a variance.
\end{proof}

\begin{lemma}[The geometry is a sphere]\label{lem:sphere}
The Fisher--Rao metric $g_{ij}=\delta_{ij}/p_i$ on the interior of $\Delta^{m-1}$ is isometric, via
$\Phi(p)=2\sqrt p$, to the positive orthant of the sphere $S^{m-1}(2)$; its constant sectional
curvature is $\tfrac14$. \forcedtag
\end{lemma}
\begin{proof}
With $\Phi(p)=(2\sqrt{p_1},\dots,2\sqrt{p_m})$ one has $d(2\sqrt{p_i})=p_i^{-1/2}dp_i$, so the
pullback of the Euclidean metric is $\sum_i(d\,2\sqrt{p_i})^2=\sum_i p_i^{-1}dp_i^2$, which is exactly
$g$. The image satisfies $\sum_i(2\sqrt{p_i})^2=4\sum_i p_i=4$, hence lies on the sphere of radius
$2$; a round sphere of radius $r$ has constant sectional curvature $1/r^2$, here $\tfrac14$
(Rao~\cite{rao}; the $m=3$ case is checked symbolically in Appendix~\ref{app:verify}). The Gibbs
curve $\beta\mapsto p(\beta)$ thus lives on a fixed curvature--$\tfrac14$ sphere, and its length
element is $\sqrt{I(\beta)}=\sqrt{\Var_\beta(\log\Mah)}$.
\end{proof}

\section{Dual flatness}\label{sec:dual}
The family is a dually flat manifold in the sense of Amari--Nagaoka \cite{amari}: the natural
($e$--affine) coordinate is $\beta$, the expectation ($m$--affine) coordinate is
$\eta=A'(\beta)=-U(\beta)$ with $U(\beta):=\E_\beta[\log\Mah]$ the mean cost, and the two dual
potentials are $A(\beta)=\log Z$ and its Legendre transform, the negentropy $A^\ast(\eta)=-H(p)$.

\begin{proposition}[Bregman form of $\mathrm{KL}$]\label{prop:bregman}
The Kullback--Leibler divergence between two temperatures is the Bregman divergence of $\log Z$:
\[
\mathrm{KL}\!\big(p(\beta_1)\,\|\,p(\beta_2)\big)=A(\beta_2)-A(\beta_1)+(\beta_2-\beta_1)\,U(\beta_1).
\]
In particular the Gibbs curve $\beta\mapsto p(\beta)$ is an $e$--geodesic (a straight line in the
natural coordinate), and $A^\ast(\eta)=-H$ holds identically. \forcedtag
\end{proposition}
\begin{proof}
For any $i$,
\[
\log\frac{p_i(\beta_1)}{p_i(\beta_2)}
=\big(-\beta_1\log\Mah_i-A(\beta_1)\big)-\big(-\beta_2\log\Mah_i-A(\beta_2)\big)
=(\beta_2-\beta_1)\log\Mah_i+A(\beta_2)-A(\beta_1).
\]
Averaging over $p(\beta_1)$ and using $\E_{\beta_1}[\log\Mah]=U(\beta_1)$,
\[
\mathrm{KL}\big(p(\beta_1)\|p(\beta_2)\big)
=(\beta_2-\beta_1)U(\beta_1)+A(\beta_2)-A(\beta_1),
\]
which is the Bregman divergence generated by the convex potential $A$. The Legendre identity
$A^\ast(\eta)=\beta\eta-A(\beta)=\sum_i p_i\log p_i=-H$ is the standard exponential--family duality.
Both are verified in Appendix~\ref{app:verify}.
\end{proof}

As a worked instance, taking $\beta_1=-1,\ \beta_2=1$ and reading $A,U$ from the thermodynamic table
of \S\ref{sec:thermo}, $\mathrm{KL}(p(-1)\|p(1))=A(1)-A(-1)+2\,U(-1)=0.950-3.256+2(1.458)=0.610$
nats---a description--length gap of $0.610$ nats between the anchor and the MDL prior.

Consequently the generalized Pythagorean theorem holds: if the $e$--geodesic from $p$ to $q$ meets
the $m$--geodesic from $q$ to $r$ Fisher--orthogonally, then
$\mathrm{KL}(p\|r)=\mathrm{KL}(p\|q)+\mathrm{KL}(q\|r)$.

\section{Thermodynamics of the catalog}\label{sec:thermo}
With energies $E_i=\log\Mah_i$ and temperature $T=1/\beta$: free energy $F(\beta)=-A(\beta)/\beta$,
internal energy $U(\beta)=-A'(\beta)$, and Gibbs entropy $S(\beta)=A(\beta)+\beta U(\beta)$, which
coincides with the Shannon entropy $H(p(\beta))$.

\begin{proposition}[Fisher information is a heat capacity]\label{prop:heat}
\[
C(\beta)=\frac{dU}{dT}=\beta^2\,I(\beta)=\beta^2\,\Var_\beta(\log\Mah)\ \ge 0 .
\]
\forcedtag
\end{proposition}
\begin{proof}
$dU/d\beta=-A''(\beta)=-I(\beta)$ and $dT/d\beta=-1/\beta^2$, so
$C=(dU/d\beta)/(dT/d\beta)=\beta^2 I(\beta)$; the non--negativity is $I\ge0$ (verified in
Appendix~\ref{app:verify}).
\end{proof}

The fluctuation $\Var_\beta(\log\Mah)$ peaks at a distinguished \emph{peak--fluctuation temperature}
$\beta^\ast\approx-0.0768$, where $\Var_{\beta^\ast}(\log\Mah)\approx0.3289$ \computedtag; the heat
capacity $C=\beta^2 I$ is however suppressed near $\beta=0$ by the $\beta^2$ factor and grows with
$|\beta|$. As $\beta\to+\infty$ the measure concentrates on the minimal--Mahler seeds $\ph,\tau$, and
as $\beta\to-\infty$ on the maximal--Mahler seed $\ph^4$. Table~\ref{tab:thermo} records the
functions at the distinguished temperatures.

\begin{table}[h]\centering
\begin{tabular}{lcccccc}
\toprule
$\beta$ & $Z(\beta)$ & $A=\log Z$ & $U=\langle\log\Mah\rangle$ & $S=H$ & $C=\beta^2 I$ & $F=-A/\beta$\\ \midrule
$-1$ & $25.944$ & $3.256$ & $1.458$ & $1.798$ & $0.2626$ & $3.256$\\
$\beta^\ast\!\approx\!-0.077$ & $7.654$ & $2.035$ & $1.176$ & $1.945$ & $0.0019$ & $26.51$\\
$0$ & $7$ & $1.946$ & $1.151$ & $1.946$ & $0$ & ---\\
$1$ & $2.586$ & $0.950$ & $0.855$ & $1.806$ & $0.243$ & $-0.950$\\
$\sqrt5$ & $1.040$ & $0.039$ & $0.648$ & $1.487$ & $0.520$ & $-0.018$\\
\bottomrule
\end{tabular}
\caption{Thermodynamics of the Mahler--Gibbs family at the distinguished temperatures ($A,U,S,C$ in
nats). Entropy $S=H$ is maximal, $\log 7\approx1.946$, at $\beta=0$; the fluctuation peaks at
$\beta^\ast$; the heat capacity $C=\beta^2 I$ vanishes at $\beta=0$ and grows on either side. All
values \computedtag.}\label{tab:thermo}
\end{table}

\section{Distances and information length}
\begin{proposition}[Fisher--Rao distance]\label{prop:dist}
For distributions $p,q$, the geodesic distance is the Bhattacharyya angle
$d_{\mathrm{FR}}(p,q)=2\arccos\sum_i\sqrt{p_i q_i}$. The seven pure states (vertices) are mutually
equidistant, $d_{\mathrm{FR}}(\delta_i,\delta_j)=\pi$ for $i\neq j$, forming a regular spherical
simplex. \forcedtag
\end{proposition}
\begin{proof}
Under $\Phi(p)=2\sqrt p$ (Lemma~\ref{lem:sphere}) the images $\Phi(p),\Phi(q)$ lie on the radius--2
sphere with Euclidean inner product $\Phi(p)\cdot\Phi(q)=4\sum_i\sqrt{p_iq_i}$; on a sphere of radius
$r$ the geodesic distance is $r$ times the central angle, here
$2\arccos\!\big(\tfrac14\Phi(p)\cdot\Phi(q)\big)=2\arccos\sum_i\sqrt{p_iq_i}$. For distinct vertices
the Bhattacharyya coefficient $\sum_k\sqrt{\delta_i(k)\delta_j(k)}=0$, so
$d_{\mathrm{FR}}=2\arccos 0=\pi$.
\end{proof}
The information length of the Gibbs curve is $L(\beta_0,\beta_1)=\int_{\beta_0}^{\beta_1}
\sqrt{I(\beta)}\,d\beta=\int\!\sqrt{\Var_\beta(\log\Mah)}\,d\beta$. Between the exact anchor and the
exchange--rate temperature, $L(-1\to\sqrt5)\approx1.5999$ nats, while the chord distance
$d_{\mathrm{FR}}(p(-1),p(\sqrt5))\approx1.5812\le L$ \computedtag, as an arc bounds its chord.
Table~\ref{tab:dist} gives the full distance matrix among the distinguished temperatures.

\begin{table}[h]\centering
\begin{tabular}{lcccc}
\toprule
$d_{\mathrm{FR}}$ & $\beta=-1$ & $\beta=0$ & $\beta=1$ & $\beta=\sqrt5$\\ \midrule
$\beta=-1$ & $0$ & $0.553$ & $1.092$ & $1.581$\\
$\beta=0$ & $0.553$ & $0$ & $0.542$ & $1.040$\\
$\beta=1$ & $1.092$ & $0.542$ & $0$ & $0.502$\\
$\beta=\sqrt5$ & $1.581$ & $1.040$ & $0.502$ & $0$\\
\bottomrule
\end{tabular}
\caption{Fisher--Rao distances $d_{\mathrm{FR}}=2\arccos\sum_i\sqrt{p_iq_i}$ (nats) between the Gibbs
distributions at the distinguished temperatures. Distances are monotone along the ordered curve but
\emph{sub--additive}: $d(-1,\sqrt5)=1.581<0.553+0.542+0.502=1.597$. This is the signature that the
Gibbs curve is an $e$--geodesic (straight in $\beta$, \S\ref{sec:dual}) but \emph{not} a Fisher--Rao
metric geodesic---the metric geodesic between the endpoints is shorter, and the full arc length
$1.600$ (the information length) exceeds the chord $1.581$. Every pure--state pair sits at the
maximal distance $\pi$. All values \computedtag.}\label{tab:dist}
\end{table}

\section{The exact anchor ($\beta=-1$, $p\propto\Mah$)}
Here $p_i=\Mah_i/Z$, $Z=17+4\sqrt5$, all algebraic in $\Q(\sqrt5)$, and the metric is the exact
diagonal $g=\operatorname{diag}(Z/\Mah_i)$.
\begin{table}[h]\centering
\begin{tabular}{lccc}
\toprule
seed & $p_i=\Mah_i/Z$ & $p_i$ & $S_i=-\log p_i$\\ \midrule
$\sqrt2$ & $2/(17+4\sqrt5)$ & $0.0771$ & $2.5628$\\
$\sqrt3$ & $3/(17+4\sqrt5)$ & $0.1156$ & $2.1573$\\
$\sqrt5$ & $5/(17+4\sqrt5)$ & $0.1927$ & $1.6465$\\
$\ph$ & $\ph/(17+4\sqrt5)$ & $0.0624$ & $2.7747$\\
$\tau$ & $\ph/(17+4\sqrt5)$ & $0.0624$ & $2.7747$\\
$\ph^4$ & $\ph^4/(17+4\sqrt5)$ & $0.2642$ & $1.3311$\\
$K$ & $(\ph^4-1)/(17+4\sqrt5)$ & $0.2256$ & $1.4888$\\
\bottomrule
\end{tabular}
\caption{The exact anchor distribution ($\beta=-1$, $p\propto\Mah$): weights exact in $\Q(\sqrt5)$,
surprisals for display. The measure favours the high--Mahler seeds $\ph^4,K$ and most suppresses the
golden pair $\ph,\tau$ (equal weight, as $\Mah(\ph)=\Mah(\tau)$). The seven weights sum to $1$.}
\label{tab:anchor}
\end{table}

\section{The keystone--power sub--model: entry and order}
With $R=\bigl(\begin{smallmatrix}0&1\\1&1\end{smallmatrix}\bigr)$, $R^n=F_nR+F_{n-1}I$ and
$\Mah(\mathrm{charpoly}\,R^n)=\ph^n$, so on the ordered sequence $\{R^n:n_0\le n\le N\}$ the cost is
linear, $E_n=n\log\ph$.
\begin{proposition}[Geometric sub--model]\label{prop:geo}
With $r=\ph^{-\beta}\in(0,1)$, the Gibbs law is the shifted geometric $p_n=(1-r)r^{\,n-n_0}$,
$Z=\dfrac{r^{n_0}-r^{N+1}}{1-r}\to\dfrac{r^{n_0}}{1-r}$, the surprisal is linear
$S_n=(n-n_0)\beta\log\ph-\log(1-r)$, and $I(\beta)=(\log\ph)^2\,r/(1-r)^2$. The point of entry $n_0$
shifts the law; the order $n$ is its support. \forcedtag
\end{proposition}

The sequence data are exact integers and their forced measures (Table~\ref{tab:keystone}); the row
$n=4$ carries the catalog gap seed $x^2-7x+1$, whose trace $L_4=7$ is the normalisation appearing
throughout the corpus.

\begin{table}[h]\centering
\begin{tabular}{ccccc}
\toprule
$n$ & $F_n$ & $L_n=\ph^n+\psi^n$ & $\Mah(R^n)=\ph^n$ & cost $E_n=n\log\ph$\\ \midrule
$1$ & $1$ & $1$ & $1.618$ & $0.481$\\
$2$ & $1$ & $3$ & $2.618$ & $0.962$\\
$3$ & $2$ & $4$ & $4.236$ & $1.444$\\
$4$ & $3$ & $7$ & $6.854$ & $1.925$\\
$5$ & $5$ & $11$ & $11.090$ & $2.406$\\
$6$ & $8$ & $18$ & $17.944$ & $2.887$\\
\bottomrule
\end{tabular}
\caption{The keystone--power sub--catalog $\{R^n\}$: Fibonacci $F_n$, Lucas $L_n$, the forced Mahler
measure $\ph^n$, and the linear cost $n\log\ph$. Linearity of the cost is what makes the Gibbs law
exactly geometric (Prop.~\ref{prop:geo}).}\label{tab:keystone}
\end{table}

\section{Two--statistic field geometry: a constant--curvature ruled surface}
Adjoin the degree $d_i$ as a second forced statistic; note $d_i=2+2\cdot\mathbf 1_{\{i=K\}}$, so
degree is affinely the indicator of the single outcome $K$. The two--parameter family
$p_i\propto\exp(-\beta_1\log\Mah_i-\beta_2 d_i)$ is an exponential family whose Fisher matrix is the
Hessian $\nabla^2 A(\beta_1,\beta_2)=\Cov_\beta(\log\Mah,\,\mathrm{deg})$. At the uniform point,
\[
g=\begin{pmatrix}0.3283 & 0.1761\\[1pt] 0.1761 & 0.4898\end{pmatrix},\qquad
\Var(\mathrm{deg})=\tfrac{24}{49}=0.4898,
\]
symmetric positive definite. This is the genuinely two--dimensional field metric; we now determine
its curvature exactly.

\begin{theorem}[Indicator--suspension]\label{thm:susp}
Let $(a,c)$ be statistics on a finite set $\Omega$ with positive base measure $\nu$, and suppose $c$
is affinely the indicator of a single outcome $K$ (i.e. $c=\alpha+\beta\,\mathbf 1_K$, $\beta\neq0$)
while $a$ is non--constant on $\Omega\setminus\{K\}$. Then the Fisher--Rao geometry of
$p_i\propto\nu_i e^{\theta_1 a_i+\theta_2 c_i}$ has \emph{constant} Gaussian curvature $\tfrac14$; it
is an isometric round $\tfrac14$--sphere, embedded in $S^{|\Omega|-1}(2)$ as a \emph{ruled} surface
swept by great circles (the $c$--direction flow lines). It is \emph{not} totally geodesic unless $a$
is affinely constant on $\Omega\setminus\{K\}$ (which fails whenever there are three or more distinct
$a$--values). \forcedtag
\end{theorem}
\begin{proof}
The Fisher metric of an exponential family depends on the statistics only through the affine span
$\langle 1,a,c\rangle$ and is invariant, as a Riemannian manifold, under affine reparametrization; so
replace $c$ by $\mathbf 1_K$ and parametrize by $(\theta,q)$, with $q=p_K\in(0,1)$ the apex mass and
$\theta$ the natural parameter of the conditional family $\rho(\theta)$ on $\Omega\setminus\{K\}$, so
$p=((1-q)\rho(\theta),q)$. Pass to square--root coordinates: $p\mapsto\sqrt p$ is an isometry onto the
unit sphere whose round metric is $\tfrac14$ the Fisher metric. With $x(\theta)=\sqrt{\rho(\theta)}$ a
unit vector ($|x|^2=\sum_i\rho_i=1$) tracing a curve on the equatorial sphere, and $\sqrt q=\sin\psi$
($\psi\in(0,\tfrac\pi2)$),
\[
\sqrt p=\big(\cos\psi\,x(\theta),\ \sin\psi\big),\qquad
\partial_\psi\sqrt p=(-\sin\psi\,x,\cos\psi),\qquad
\partial_\theta\sqrt p=(\cos\psi\,x',0).
\]
Using $|x|^2=1\Rightarrow x\cdot x'=0$, the induced unit--sphere metric is
\[
g^{\text{unit}}=\underbrace{(\sin^2\!\psi\,|x|^2+\cos^2\!\psi)}_{=\,1}\,d\psi^2
+2\underbrace{(-\sin\psi\cos\psi\,x\!\cdot\! x')}_{=\,0}\,d\psi\,d\theta
+\cos^2\!\psi\,|x'|^2\,d\theta^2
= d\psi^2+\cos^2\!\psi\,d\ell^2,
\]
reparametrizing $\theta$ by the arc length $\ell$ of $x$ ($d\ell=|x'|\,d\theta$). This is the round
unit--sphere metric (latitude $\psi$, longitude $\ell$), of Gaussian curvature $-f''/f=1$ with
$f=\cos\psi$---\emph{independent of the curve} $x$. As the Fisher metric is $4\,g^{\text{unit}}$, its
curvature is the constant $\tfrac14$. The family is thus the spherical suspension of the apex $K$ over
the $a$--curve; non--constancy of $a$ makes $x(\theta)$ a genuine immersed curve, so the surface is
two--dimensional.
\end{proof}

The embedding structure is exactly a ruled surface. Writing $s=\sqrt p$ and $\tau_\alpha=T_\alpha-\langle
T_\alpha\rangle$ (centred statistics), the second fundamental form in the sphere is $\mathrm{II}_{\alpha
\beta}=\tfrac14(\tau_\alpha\tau_\beta\,s)^\perp$, and $\mathrm{II}_{\alpha\beta}=0$ iff the pointwise
product $T_\alpha T_\beta\in V:=\langle 1,a,c\rangle$. For $c=\mathbf 1_K$: $c^2=c\in V$ (idempotent)
and $a\,c=a_K\,c\in V$ (the indicator localises the product), so $\mathrm{II}_{cc}=\mathrm{II}_{ac}=0$
---the $c$--flow lines are great circles and the surface is ruled by them. But $a^2\notin V$ whenever
$a$ has $\ge3$ distinct values (Prop.~\ref{prop:notTG}), so $\mathrm{II}_{aa}\neq0$: the surface bends
in the $a$--direction and is not totally geodesic. This is why the curvature is constant $\tfrac14$
(the extrinsic curvature $\det\mathrm{II}=\langle\mathrm{II}_{aa},\mathrm{II}_{cc}\rangle-|\mathrm{II}
_{ac}|^2=0$) without the surface being a great subsphere.

\begin{corollary}[the catalog surface]\label{cor:catsurf}
Since $\mathrm{deg}=2+2\cdot\mathbf 1_{\{i=K\}}$ (only the quartic $K$, minimal polynomial
$x^4+5x^2-5$, has degree $4$) and $\log\Mah$ is non--constant on the remaining seeds, the catalog's
$(\log\Mah,\mathrm{deg})$ Gibbs surface satisfies Theorem~\ref{thm:susp} and has \emph{exactly}
constant curvature $\tfrac14$. The multiplicity of $\ph,\tau$ (equal $\log\Mah$ and degree) is
absorbed as a base--measure weight and does not affect the result. \forcedtag
\end{corollary}

\paragraph{The catalog surface, explicitly.} At the anchor $\beta=-1$ the apex mass is
$q_K=p_K=0.2256$, giving latitude $\psi=\arcsin\sqrt{q_K}=0.495$; the conditional distribution on the
six non--apex seeds is $(0.0996,0.1493,0.2489,0.0805,0.0805,0.3412)$, whose square roots are the
point $x$ on the equatorial $S^5$. One checks $\cos^2\psi=1-q_K=0.7744$, so the catalog point sits on
the suspension exactly as $\sqrt p=(\cos\psi\,x,\sin\psi)$. Sweeping $\beta$ moves $x$ along the
$\log\Mah$--curve and $\psi$ toward the apex $K$, tracing out the constant--curvature--$\tfrac14$
sphere; the coordinate $(\psi,\ell)$ of Theorem~\ref{thm:susp}, with $\ell$ the arc length
of the $\log\Mah$--curve, realises it as a round sphere intrinsically (though not as a great subsphere
of the ambient).

This inverts the naive expectation: the second statistic does not bend the geometry into a new
varying--curvature object; the $(\log\Mah,\mathrm{deg})$ family sits inside the surprisal sphere as a
constant--curvature ruled surface. The apex is essential---a generic two--statistic family (no
indicator) has non--constant, sign--changing curvature; explicit controls give $4/21$ ($m=4$) and
$857/16928$ ($m=5$), both $\neq\tfrac14$. The general Brioschi curvature of an indicator family, its
invariance under base measure (hence constancy), the ruled second fundamental form, and the suspension
metric are all machine--verified (\texttt{suspension\_theorem.py} and
\texttt{field\_surprisal\_classification.py}).

\begin{remark}
Neither the trace form nor the $(\log\Mah,\mathrm{deg})$ surface is totally geodesic in the surprisal
sphere---but for opposite reasons. Totally geodesic surfaces are the curvature--$\tfrac14$ subspheres;
the field's \emph{flat} trace--form geometry misses this because its curvature is $0\neq\tfrac14$
(forcing extrinsic $k_1k_2=-\tfrac14$, \S12), while the $(\log\Mah,\mathrm{deg})$ surface has the right
intrinsic curvature $\tfrac14$ but the wrong embedding---it is ruled, not geodesic ($\mathrm{II}_{aa}
\neq0$). In fact (Prop.~\ref{prop:notTG}) \emph{no} $(\log\Mah,X)$ surface is totally geodesic.
\end{remark}

\section{The multi-statistic landscape: degree is geometrically distinguished}\label{sec:landscape}
Theorem~\ref{thm:susp} singles out one property of the degree statistic---being an indicator of a
\emph{single} outcome. The catalog carries other forced invariants, and comparing their two--statistic
surfaces shows how special that property is (Table~\ref{tab:landscape}).

\begin{table}[h]\centering
\begin{tabular}{llll}
\toprule
statistic $X$ & values & type & $K$ of $(\log\Mah,X)$ at three points\\ \midrule
degree & $\{2,4\}$ & $\mathbf 1_{\{K\}}$ (single outcome) & $0.250,\ 0.250,\ 0.250$\\
$\#\{|\zeta|>1\}$ & $\{1,2\}$ & $\mathbf 1_{\{\ph,\tau,\ph^4\}}$ ($3$--set) & $0.020,\ 0.066,\ -0.006$\\
trace & $\{-1,0,1,7\}$ & non--indicator & $0.177,\ 0.056,\ 0.197$\\
house $\max|\zeta|$ & six values & non--indicator & $0.211,\ 0.114,\ 0.227$\\
\bottomrule
\end{tabular}
\caption{Gaussian curvature of the two--statistic Gibbs surface $(\log\Mah,X)$ for four forced
invariants, sampled at three natural--parameter points. Only the single--outcome indicator (degree)
is the constant $\tfrac14$. All \computedtag{} (curvature routine validated on four metrics of known
curvature).}\label{tab:landscape}
\end{table}

The number of conjugates outside the unit circle is itself an indicator, $\#\{|\zeta|>1\}
=2-\mathbf 1_{\{\ph,\tau,\ph^4\}}$---but of a \emph{three--element} set, not a single outcome. Its
surface has non--constant, sign--changing curvature, so the single--outcome hypothesis of
Theorem~\ref{thm:susp} is sharp: a set indicator makes the family a join of two non--trivial curves
(one on the set's sub--sphere, one on its complement), which is not a suspension over a point, hence
not a round sphere. The trace and house are generic (non--indicator) and likewise vary. Among the
catalog's forced invariants, \emph{degree is the unique one whose Gibbs surface with $\log\Mah$ has
constant curvature $\tfrac14$}---a geometric distinction of the single quartic seed $K$, verified in
\texttt{field\_surprisal\_tier2.py} ($15/15$).

\section{The charge grading: a second constant--curvature--$\tfrac14$ surface}\label{sec:charge}
The catalog carries a further forced invariant from the companion charge--measure coupling
\cite{cmc}: the conjugate--angle charge. An algebraic number is charge--admissible if every conjugate
$\alpha$ satisfies $\alpha^n\in\R_{>0}$ for some $n\ge1$; the least such $n$ is its charge group
$\Z/n\Z$, and the charge of a root at argument $\theta$ is $\chi_n=\operatorname{round}(n\theta/2\pi)
\bmod n$. Read off the catalog:
\[
\begin{array}{lccccccc}
\text{seed} & \sqrt2 & \sqrt3 & \sqrt5 & \ph & \tau & \ph^4 & K\\[1pt]
\text{charge group} & \Z/2 & \Z/2 & \Z/2 & \Z/2 & \Z/2 & \Z/1 & \Z/4
\end{array}
\]
The seed $\ph^4$ is the \emph{unique} one with both conjugates real positive (roots $\ph^4,\ph^{-4}$),
so its charge group is trivial $\Z/1$; $K$ alone has complex conjugates $\pm i\beta$, giving the
maximal $\Z/4$; the remaining five sit at $\Z/2$.

The paper's central grading is charge \emph{parity}, which sets the Mahler floor---$\ph$ for even $n$,
$2$ for odd $n$, the odd gap $(\ph,2)$ being the no--Salem closure \cite[Thm.~5.1, 6.2]{cmc}. For this
catalog parity is $\mathbf 1_{\{\ph^4\}}$, a single--outcome indicator, so by Theorem~\ref{thm:susp}
the surface $(\log\Mah,\text{charge parity})$ again has constant curvature $\tfrac14$---verified
at $K=0.250,0.250,0.250$, and exactly, the Brioschi curvature of an arbitrary statistic with the
$\ph^4$ indicator being identically $\tfrac14$ \forcedtag. The charge group \emph{order}
$\{1,2,4\}$ is three--valued and gives a non--$\tfrac14$, varying surface ($0.339,0.362,0.342$),
consistent with \S\ref{sec:landscape}.

We can now classify \emph{all} such surfaces completely, and the answer corrects a natural
over--expectation: constant curvature $\tfrac14$ is strictly weaker than being totally geodesic (a great
subsphere), and here only the former ever holds.

\begin{proposition}[No surface is totally geodesic]\label{prop:notTG}
For every second statistic $X$, the $(\log\Mah,X)$ surface is \emph{not} totally geodesic in
$S^{6}(2)$. Indeed a totally geodesic $V=\langle 1,\log\Mah,X\rangle$ would be a unital subalgebra of $\R^7$,
forcing the number of distinct value pairs $(\log\Mah_i,X_i)$ down to $\dim V=3$; but the cost takes
six distinct values, so the pair count is always $6$ or $7$, never the $3$ a great $2$--subsphere
requires. \forcedtag
\end{proposition}

\begin{remark}[Correction]\label{rem:corr}
An earlier draft argued $\mathrm{II}_{aa}\neq0$ from ``$\log^2\Mah\notin V$''; that inference is false
as stated: for $X$ affinely $\log^2\Mah$ one has $\log^2\Mah\in V$ and $\mathrm{II}_{aa}=0$ (then
$\mathrm{II}_{ac}\neq0$ instead, since $\log^3\Mah\notin\langle1,\log\Mah,\log^2\Mah\rangle$;
machine--checked in \texttt{t1\_core.py}). The value--count argument above is the correct and complete
proof; the proposition's statement is unchanged.
\end{remark}

\begin{theorem}[Complete constant--curvature classification]\label{thm:classification}
The $(\log\Mah,X)$ surface has constant Gaussian curvature $\tfrac14$ iff $X$ is affinely the indicator
$\mathbf 1_S$ of a set $S$ with $\log\Mah$ constant on $S$ or on its complement. There are
\emph{exactly eight} such surfaces (a surface depends only on the partition $\{S,S^c\}$): the seven
single--seed indicators and the golden--pair merge $\mathbf 1_{\{\ph,\tau\}}$. Each is a ruled
constant--curvature--$\tfrac14$ surface (Theorem~\ref{thm:susp}), none totally geodesic. No
\emph{non}--indicator $X$ gives constant curvature $\tfrac14$.
\end{theorem}
\noindent\emph{Proof.} Sufficiency is Theorem~\ref{thm:susp} (idempotency gives
$\mathrm{II}_{cc}=0$; $\log\Mah$ constant on $S$ gives $\mathrm{II}_{ac}=0$; hence $\det\mathrm{II}=0$);
the $16$--subset census among all $127$ indicators is exhaustive (\texttt{field\_surprisal\_classification.py},
$9/9$). Necessity---no non--indicator $X$ is constant--$\tfrac14$---is proved in full in
\S\ref{sec:classdone} (Theorem~\ref{thm:classfull}); the earlier $\approx5000$--sample evidence is
retired and the tag upgrades to \forcedtag.

\begin{remark}[The eight surfaces and their forcing]
Seven of the eight are single--seed, and each is \emph{forced} by an arithmetic invariant:
\[
\begin{array}{lccccccc}
\text{seed } s & \sqrt2 & \sqrt3 & \sqrt5 & \ph & \tau & \ph^4 & K\\[1pt]
\text{singled out by} & \Mah{=}2 & \Mah{=}3 & \Mah{=}5 & \mathrm{tr}{=}1 & \mathrm{tr}{=}{-}1 & \Mah{=}\ph^4 & \Mah{=}\ph^4{-}1
\end{array}
\]
The Mahler measure alone distinguishes five seeds; the only degeneracy is the golden pair
$\{\ph,\tau\}$ tied at $\Mah=\ph$, split by the trace $\pm1$---the same sign the $\lambda=2c$
construction uses to select the Perron keystone $R^2=R+I$ (trace $+1$) over the Clifford conjugate gate
$R^2=I-R$ (trace $-1$) \cite{lambda}. The eighth surface is precisely $\mathbf 1_{\{\ph,\tau\}}$, the
single--outcome indicator of the \emph{effective} catalog in which the Mahler--tied seeds $\ph,\tau$ are
merged. The two deepest of the seven are $K$ and $\ph^4$, the charge--extreme seeds (charge group $\Z/4$
and $\Z/1$), additionally distinguished by degree, non--reality, and total positivity.
\end{remark}

\section{The classification completed: a finite branch calculus}\label{sec:classdone}
Theorem~\ref{thm:classification} left one \opentag: ruling out every non--indicator statistic. We
close it. Write $a=\log\Mah$, $V=\langle1,a,X\rangle$ (assumed $3$--dimensional), and let $Q_\theta$
be the $L^2(p_\theta)$--projection onto $V^\perp$ along the family $p\propto e^{\theta_1a+\theta_2X}$.
The Gauss equation reduces constancy to one scalar: with $F=a^2$, $G=X^2$, $H=aX$,
\[
q(\theta)\;=\;\langle Q_\theta F,\,Q_\theta G\rangle-\|Q_\theta H\|^2,
\qquad K\equiv\tfrac14\iff q\equiv0 .
\]

\begin{lemma}[$\theta$--rigidity]\label{lem:rigid}
$\ker Q_\theta=V$ for every $\theta$ (the Gram matrix of $(1,a,X)$ is nonsingular under any positive
weight). Hence $\mathrm{II}_{ac}\equiv0\iff aX\in V$ and $\mathrm{II}_{cc}\equiv0\iff X^2\in V$: each
ruled condition holds at one temperature iff at all. \forcedtag
\end{lemma}

\begin{theorem}[The ruled system, solved]\label{thm:ruled}
$\{aX\in V,\;X^2\in V\}\iff X\!\cdot\!V\subseteq V\iff X$ is affinely $\mathbf 1_S$ with $a$ constant
on $S$ or $S^c$---exactly the eight families. Module route: $X\!\cdot\!V\subseteq V$ makes $V$ a module
over itself, hence a sum of level--set blocks of $X$; $\dim V=3$ against six distinct costs leaves
two levels with $a$ constant on one. Elimination route: $aX\in V$ makes $X$ a M\"obius function
$\mu(a)$ off at most one pole level; the increment identity
$\mu(s)-\mu(t)=(d_0+d_1d_2)(t-s)/\{(s-d_2)(t-d_2)\}$ and $X^2\in V$ force $\mu$ constant off--pole,
leaving pole--level indicators. \forcedtag\ (\texttt{t1\_core.py}: the invariant module
$N(X)=\{v\in V:vV\subseteq V\}$ is $2$--dimensional exactly on the eight families, $1$--dimensional on
every control, including $X\simeq\log^2\Mah$ and the pair--pole family.)
\end{theorem}

Off the ruled locus, $q\equiv0$ must hold with both forms nonzero. Two exact reductions make the
question finite.

\begin{proposition}[Numerator and the window identity]\label{prop:window}
Put $C(f,g)=\sum_{i<j}w_iw_j(f_i-f_j)(g_i-g_j)$ on weights $w_i=e^{\theta_1a_i+\theta_2X_i}$,
$D=C(a,a)C(X,X)-C(a,X)^2$, and $R(f,g)=D\,C(f,g)-C(f,\cdot)\,\mathrm{adj}\,C\;C(\cdot,g)$. Then
$Z^2Dq=P(w):=R(F,G)-R(H,H)$, a degree--$6$ form whose $462$ monomials all have support $\ge4$. For a
four--set $s$ let $\ell_s$ be the signed $3\times3$ minors of the rows $(1,a,X)|_s$ and
\[
q_4(s)\;=\;\ell_s(a^2)\,\ell_s(X^2)-\ell_s(aX)^2 .
\]
Then the coefficient of $w_h^3w_jw_kw_l$ in $P$ \emph{equals} $q_4(s)$, cofactor exactly $1$, for all
$140$ heavy placements; and on the four--point face itself $\langle Qf,Qg\rangle=\kappa\,
\ell_s(f)\ell_s(g)$ with $\kappa>0$, so $q=\kappa\,q_4(s)$ there. \forcedtag\
(\texttt{t1\_engine.py}, \texttt{t1\_reduction.py}.)
\end{proposition}

\begin{theorem}[Master window identity: the coefficient law is Sylvester plus Cauchy--Binet]\label{thm:master}
As a polynomial identity in independent weights $w_1,\dots,w_m$ and arbitrary statistics $(a,X)$,
\[
P(w)\;=\;Z^2\sum_{|s|=4}q_4(s)\,w_s,\qquad Z=\textstyle\sum_iw_i,\ \ w_s=\textstyle\prod_{i\in s}w_i,\qquad
q_4(s)=\ell_s(a^2)\,\ell_s(X^2)-\ell_s(aX)^2,
\]
where $\ell_s(f)=\det\,(\mathbf 1,a,X,f)|_s$ realizes the signed--minor functional of
Proposition~\ref{prop:window}. \forcedtag
\end{theorem}
\begin{proof}
Three classical steps. \emph{(i)} $R(f,g)$ is the $3\times3$ determinant of the matrix
$\bigl(C(u,v)\bigr)_{u\in(f,a,X),\,v\in(g,a,X)}$: expanding along the first row reproduces
$D\,C(f,g)-C(f,\cdot)\,\mathrm{adj}\,C\,C(\cdot,g)$ term by term. \emph{(ii)} With
$\langle u,v\rangle_w=\sum_iw_iu_iv_i$ one has
$C(u,v)=Z\langle u,v\rangle_w-\langle u,\mathbf 1\rangle_w\langle\mathbf 1,v\rangle_w$, the
$2\times2$ minor of the mixed Gram matrix
$B(f,g)=\bigl(\langle u,v\rangle_w\bigr)_{u\in(\mathbf 1,f,a,X),\,v\in(\mathbf 1,g,a,X)}$ bordered on
its $(1,1)$ entry $\langle\mathbf 1,\mathbf 1\rangle_w=Z$. Sylvester's determinant identity with
pivot $Z$ \cite{hj} gives $\det\bigl(C(u,v)\bigr)=Z^{4-2}\det B(f,g)$ (reordering rows and columns
simultaneously leaves determinants unchanged). \emph{(iii)} Writing $B(f,g)=UWV^{\mathsf T}$ with
$U,V$ the $4\times m$ statistic matrices and $W=\operatorname{diag}(w)$, Cauchy--Binet \cite{hj}
gives $\det B(f,g)=\sum_{|s|=4}w_s\,\ell_s(f)\,\ell_s(g)$. Hence
$P=R(F,G)-R(H,H)=Z^2\sum_sw_s\bigl[\ell_s(F)\ell_s(G)-\ell_s(H)^2\bigr]$.
\end{proof}
\begin{corollary}[Coefficient law and structure]\label{cor:windowlaw}
(a) Every monomial of $P$ has degree $6$ and support $\ge4$, with the exact law
$[w^\kappa]P=\sum_{|s|=4,\ \mathbf 1_s\le\kappa}c(\kappa-\mathbf 1_s)\,q_4(s)$, where $c(2e_i)=1$ and
$c(e_i+e_j)=2$; the monomial census $140+210+105+7=462$ over the shapes
$(3,1,1,1),(2,2,1,1),(2,1,1,1,1),(1^6)$ is forced. (b) The heavy coefficient
$[w_h^3w_jw_kw_l]\,P=q_4(s)$ with cofactor exactly $1$: $w_h^2$ is the unique contribution of $Z^2$.
(c) Since the squarefree monomials $\{w_s\}$ are linearly independent, $q\equiv0$ iff $q_4(s)=0$ for
\emph{all} $35$ windows---with no collision analysis. (d)
$q_4(s)=-\prod_{t\in s}\det\,(\mathbf 1,a,X)|_{s\setminus t}$, a product of the four triple minors,
so the four--point trichotomy (Lemma~\ref{lem:tri}) is immediate: $q_4(s)=0$ iff some three of the
four plane points are affinely dependent (collinear, or a coincident pair as at the golden merge).
(e) $D=Z\sum_{|t|=3}w_t\,\Delta_t^2$ by the same two identities one order down, so
$q=\bigl[\sum_sw_sq_4(s)\bigr]/\bigl[Z\sum_tw_t\Delta_t^2\bigr]$: the $\alpha=0$ curvature
obstruction is a ratio of consecutive Cauchy--Binet compound sums (the Hessian--metric picture of
\cite{shima}), and restricting $w$ to a four--point face recovers $q=\kappa\,q_4(s)$ with the
explicit $\kappa=w_s/\bigl(Z\sum_{t\subset s}w_t\Delta_t^2\bigr)>0$. \forcedtag
\end{corollary}
\begin{remark}[Architecture; Open Problem 3 resolved]\label{rem:master}
Theorem~\ref{thm:master} replaces the machine discovery of Proposition~\ref{prop:window} ($140$
exact divisions) by a two--line determinantal mechanism and resolves Open Problem (3) of the
previous edition. The classification chain now runs rigidity $\to$ ruled system $\to$
Theorem~\ref{thm:master} $\to$ trichotomy $\to$ plane lemma; the collision pencil and branch
survival (Proposition~\ref{prop:pencil}, Lemma~\ref{lem:survival}) are no longer load--bearing and
are retained as an independent verification lane and as results on the collision structure in their
own right. Verified end to end in \texttt{t1\_windowproof.py} ($44/44$, exit $0$): the chain
symbolically at $m=4,5$; the factorization (d) symbolically for four generic points; exact rational
instances at $m=6,7$; and, on the catalog ring, the full $462$--coefficient dictionary equality of
$P$ with $Z^2\sum_sq_4(s)\,w_s$, string--exact against the \texttt{t1\_engine.py} artifact.
\end{remark}

\begin{proposition}[Collision pencil]\label{prop:pencil}
$q\equiv0$ iff the coefficients of $P$ vanish after grouping monomials with equal frequency pairs
$(\sum_ik_ia_i,\ \sum_ik_iX_i)$. Since $\log2,\log3,\log5,\log\ph$ are $\Q$--linearly independent
(multiplicative independence by norms in $\Q(\sqrt5)$; no transcendence input), the cost collisions
form a rank--$2$ lattice generated by the golden swap $e_\ph-e_\tau$ and the Salem square
$2e_K-e_{\sqrt5}-e_{\ph^4}$, reflecting $(\ph^4-1)^2=5\ph^4$. Every $X$--collision condition lies on
the pencil $j\Delta_1+k\Delta_2=0$ with $\Delta_1=X_\ph-X_\tau$, $\Delta_2=2X_K-X_{\sqrt5}-X_{\ph^4}$;
the realizable primitive relations are only $\Delta_1,\ \Delta_2,\ \Delta_1\pm\Delta_2,\
2\Delta_1\pm\Delta_2$, and two independent relations force $\Delta_1=\Delta_2=0$. \forcedtag
\end{proposition}

\begin{lemma}[Branch survival]\label{lem:survival}
In every branch of the pencil---generic, each of the five nontrivial single relations, and both
sub--branches of the merged six--point recursion for $\Delta_1=0$ (whose own collision lattice has
rank $1$, again the Salem square; the base weight $2$ at the merged point is invisible to vanishing
since the collision support misses it)---every window $s$ retains at least one heavy placement whose
group is a singleton. Hence $q\equiv0$ forces $q_4(s)=0$ for \emph{all} windows in \emph{every}
branch. \forcedtag\ (\texttt{t1\_branches.py}: zero windows lost across all eight branch runs; all
$35$ windows vanish identically on the eight families in symbolic $(\alpha,\beta,\gamma)$, and a
random control violates one.)
\end{lemma}

\begin{lemma}[Four--point trichotomy]\label{lem:tri}
$q_4(s)=0$ iff some three of the four plane points $T_i=(a_i,X_i)$, $i\in s$, are collinear, or
$\{\ph,\tau\}\subset s$ with $T_\ph=T_\tau$. Collinear triples are automatically non--vertical, the
cost having a single tie.
\end{lemma}
\begin{proof}
$q_4$ is the determinant of the face's residual scalar form; it vanishes iff
$\operatorname{rank}(1,a,X)|_s\le2$ (all four collinear) or a null direction $v=\alpha a+\beta X$ has
$vV|_s\subseteq V|_s$. Diagonalizing $v$: four distinct $v$--levels are impossible ($\mathbf 1$ has
full support); three levels, or a $2{+}2$ split, force the tied pair constant in both coordinates,
i.e.\ $T_\ph=T_\tau$; a $1{+}3$ split puts $X$ affine in $a$ on the triple. The ray $\beta=0$ reduces
to $V|_s=\langle1,a,a^2\rangle$ and the minimal polynomial of $a|_s$, landing in the same two
alternatives.
\end{proof}

\begin{lemma}[Plane lemma]\label{lem:plane}
If $n\ge5$ plane points are such that every $4$--subset contains a collinear triple, then all but at
most one lie on a single line.
\end{lemma}
\begin{proof}
Let $L$ be a line through the most points, $|L|\ge3$. Suppose $p\ne q$ lie off $L$. For any
$x,y\in L$, the $4$--set $\{p,q,x,y\}$ contains a collinear triple; $\{p,x,y\}$ or $\{q,x,y\}$ would
put $p$ or $q$ on $L$, so $\mathrm{line}(p,q)$ contains $x$ or $y$---for every pair $x,y\in L$.
Taking three points of $L$ pairwise, $\mathrm{line}(p,q)$ would contain two of them and meet $L$
twice: contradiction.
\end{proof}

\begin{theorem}[Complete classification, closed]\label{thm:classfull}
$q\equiv0$ iff $X$ is affinely $\mathbf 1_S$ with $S$ one of the seven seeds or the golden pair.
Theorem~\ref{thm:classification} therefore holds as stated with tag \forcedtag: the eight surfaces
are the only constant--curvature--$\tfrac14$ two--statistic surfaces.
\end{theorem}
\begin{proof}
If $\Delta_1\ne0$ the coincidence alternative of Lemma~\ref{lem:tri} cannot fire, so by
Lemmas~\ref{lem:survival}--\ref{lem:plane} at least six of the seven $T_i$ lie on one non--vertical
line $\ell$, i.e.\ $X=\alpha+\beta a+\gamma\mathbf 1_S$ with $S$ the exceptional point. Both golden
points on $\ell$ would give $X_\ph=X_\tau$ (equal abscissae on a non--vertical line), contradicting
$\Delta_1\ne0$; so the exception is $\ph$ or $\tau$, automatically an $a$--level. If $\Delta_1=0$ the
family is the merged six--point catalog, whose six costs are distinct; the same
window--trichotomy--plane chain on its $15$ windows puts five of six points on a line, giving the
merged singletons---pulling back to the golden merge, or a singleton with $X_\ph=X_\tau$. All cases
land in the eight families; the converse is Theorem~\ref{thm:susp}.
\end{proof}

\begin{remark}[Below the ceiling]
For \emph{any} subset $S$, idempotency gives $\mathrm{II}_{cc}=0$, hence
$\det\mathrm{II}=-|\mathrm{II}_{ac}|^2\le0$ and $K\le\tfrac14$ pointwise: every indicator surface sits
at or below the $\tfrac14$ ceiling, the eight families being exactly the equality case---consistent
with the $3$--set indicator row of Table~\ref{tab:landscape}.
\end{remark}

\section{Product geometry: independent catalogs do not curve each other}
Two catalogs $\Omega_1,\Omega_2$ with their own temperatures give the product family
$p_{ij}(\beta_1,\beta_2)\propto\Mah^{(1)}_i{}^{-\beta_1}\Mah^{(2)}_j{}^{-\beta_2}
=p^{(1)}_i(\beta_1)\,p^{(2)}_j(\beta_2)$. Its log--partition splits, $A(\beta_1,\beta_2)=A_1(\beta_1)
+A_2(\beta_2)$, so the Fisher matrix is block--diagonal:
\[
\nabla^2 A=\begin{pmatrix}\Var_{\beta_1}(\log\Mah^{(1)}) & 0\\[2pt] 0 & \Var_{\beta_2}(\log\Mah^{(2)})\end{pmatrix},
\qquad \frac{\partial^2 A}{\partial\beta_1\,\partial\beta_2}=0 ;
\]
the two fields' description costs are statistically independent.

\begin{proposition}[No cross--curvature]\label{prop:product}
The joint categorical model of two independent catalogs is the Riemannian product
$S^{m_1-1}(2)\times S^{m_2-1}(2)$: its sectional curvature is $\tfrac14$ within each spherical factor
and $0$ on every mixed $2$--plane; equivalently, all Christoffel symbols mixing the two blocks vanish.
\forcedtag
\end{proposition}
\begin{proof}
The product metric is block--diagonal with each block a function of its own coordinates only; hence
every Christoffel symbol with indices spanning both blocks vanishes, so does every mixed Riemann
component, and the sectional curvature of a plane spanning one direction from each factor is $0$
(verified in \texttt{field\_surprisal\_tier2.py}). Each factor keeps its intrinsic curvature
$\tfrac14$ (Lemma~\ref{lem:sphere}).
\end{proof}

To make two fields interact geometrically---nonzero cross--curvature---one must replace the product by
a warped metric $g_1\oplus f(\cdot)^2 g_2$ (or a nontrivial fibration), introducing a coupling
function $f$. No invariance selects $f$: by the same \v{C}encov argument that leaves the temperature
$\beta$ (and the companion paper's conformal scale $c$) undetermined, the coupling is a \emph{declared}
datum \declaredtag, of exactly that epistemic status. The product of the two Gibbs \emph{curves} (one
temperature each) is, by contrast, flat---a product of one--dimensional manifolds---a degenerate
special case carrying no field interaction at all.

\section{The temperature is a coordinate, not a parameter}\label{sec:tempcoord}
The ledger has carried one \declaredtag{} since \S3: the Gibbs form $p\propto\Mah^{-\beta}$ and the
choice of $\beta$. Product functoriality dissolves the first half.

\begin{theorem}[Forced form]\label{thm:forcedform}
Let $p_i\propto f(\Mah_i)$ depend on the seeds through the Mahler measure alone, with
$f:\R_{>0}\to\R_{>0}$ Lebesgue measurable, and require product functoriality: on every product
catalog the joint weight matrix $[f(\Mah_i\Mah'_j)]$ factorizes as a product measure. Then
$f(x)=c\,x^{-\beta}$: the Mahler--Gibbs family is \emph{forced}, and $\beta$ is the canonical affine
coordinate of its $e$--geodesic. \forcedtag\ (modulo the stated regularity)
\end{theorem}
\begin{proof}
Factorization is the rank--$1$ condition on $[f(\Mah_i\Mah'_j)]$, i.e.\ $f(xy)=f(x)f(y)/f(1)$ on the
value grid; ranging over finite mass sets this is Cauchy's multiplicative equation on $\R_{>0}$,
whose measurable solutions are the power laws \cite{aczel}. The regularity is genuinely necessary: on
the bare multiplicative lattice generated by $\{2,3,5,\ph\}$, independent per--generator exponents are
multiplicative without being any single power (\texttt{t2\_temperature.py} exhibits one exactly), so
``a single $\beta$'' is precisely the content of measurability.
\end{proof}

What remains declared is only the selection of an operating \emph{point} on a forced curve, and the
curve's marked points are each forced by a separate principle:
\[
\begin{array}{lllr}
\beta & \text{forced by} & I(\beta) & s(\beta)\\[2pt]
-1 & \text{anchor } Z(-1)=17+4\sqrt5 & 0.26260 & -0.5539\\
\beta^\ast\approx-0.07677 & \text{max Fisher }(I'=-\kappa_3=0) & 0.32886 & -0.0440\\
0 & \text{maximum entropy} & 0.32832 & 0\\
1 & \text{MDL} & 0.24304 & +0.5430\\
\sqrt5 & \text{exchange rate }\lambda=2c & 0.10405 & +1.0461\\
\beta_C\approx2.5455 & \text{heat--capacity peak }(\beta>0) & \text{---} & +1.1401
\end{array}
\]
with $Z(+1)=\sum\Mah^{-1}=\tfrac{91}{30}-\tfrac{\sqrt5}{5}$ exact. \forcedtag{} for the anchors and
the identity $I'(\beta)=-\kappa_3(\beta)$; \computedtag{} for the numeric column.

\begin{proposition}[Finite total information length]\label{prop:ltot}
$I(\beta)\to0$ exponentially in both directions (spectral--gap concentration), so the Gibbs curve has
finite total Fisher--Rao length
\[
L_{\mathrm{tot}}=\int_{-\infty}^{\infty}\!\sqrt{I(\beta)}\,d\beta\;\approx\;5.64622 ,
\]
an intrinsic invariant of the catalog. Its endpoints are forced: $\beta\to+\infty$ concentrates on the
cost minimizers---the golden merge $\{\ph,\tau\}$, uniform---and $\beta\to-\infty$ on the cost
maximizer, the $\ph^4$ vertex. The forced $e$--geodesic runs from the eighth surface's apex to the
charge--trivial seed, and arc length $s$ is its canonical parametrization; the segment
$[-1,\sqrt5]$ carrying all six marked points uses only $28.3\%$ of the total length. \forcedtag{} for
the limits; \computedtag{} for the values (\texttt{t2\_temperature.py}).
\end{proposition}

\begin{theorem}[Selection dichotomy]\label{thm:dichotomy}
The metric selects no operating point: a finite--length one--dimensional Riemannian curve has
isometry group $\{\mathrm{id},\mathrm{flip}\}$, and the flip swaps the endpoints, which are
distinguishable (the $\beta\to+\infty$ limit is uniform on the golden pair, the
$\beta\to-\infty$ limit is the $\ph^4$ vertex); the labeled--endpoint isometry group is therefore
trivial, and no interior point is metrically canonical. \forcedtag\ Meanwhile the invariant
principles on offer provably disagree: with $s(0)=0$, the arc--length midpoint sits at
$s(\beta_{\mathrm{mid}})\approx-0.3401$, the max--Fisher point at $s(\beta^\ast)\approx-0.0440$,
maximum entropy at $0$, and the heat--capacity peak at $s(\beta_C)\approx+1.1401$: pairwise
distinct with margins far above numeric error. \computedtag\ Selecting an operating temperature is
therefore \emph{irreducibly} an election among inequivalent invariant principles; the
\declaredtag{} on the selection is not a gap awaiting closure but the honest form of the answer
(\texttt{t6\_selection.py}, $18/18$: exact $\Q(\sqrt5)$ anchors, symbolic $I'=-\kappa_3$,
dps--$30/50$ stability).
\end{theorem}

\section{Higher statistics: the iterated join and the geodesic threshold}\label{sec:join}
\begin{theorem}[Iterated join]\label{thm:join}
Let $s_1,\dots,s_{k-1}$ be distinct outcomes and $a$ non--constant on the complement. The
$k$--parameter family $p\propto\nu\exp(\theta_0a+\sum_j\theta_j\mathbf 1_{s_j})$ maps under
$\sqrt{p}$ onto the spherical join of the totally geodesic corner $S^{k-2}$ spanned by the apexes
with the base curve of the $a$--family on the complement; in join coordinates the metric is
\[
d\Psi^2+\sin^2\!\Psi\,g_{S^{k-2}}+\cos^2\!\Psi\,d\ell^2 ,
\]
of constant sectional curvature $1$. The family is therefore, isometrically, an open subset of a
round $\tfrac14$--sphere of dimension $k$: the suspension theorem iterates in every dimension.
\forcedtag\ (\texttt{t3\_suspension.py}: $R_{ijkl}=g_{ik}g_{jl}-g_{il}g_{jk}$ exactly for $k=3,4$;
the pullback of the round metric under the join map is computed symbolically.)
\end{theorem}

\begin{corollary}[The catalog three--fold]\label{cor:threefold}
$(\log\Mah,\ \deg,\ \text{charge parity})=(\mathrm{cost},\ \mathbf 1_K,\ \mathbf 1_{\ph^4})$ is a
three--dimensional round $\tfrac14$--sphere: cost, degree structure, and total positivity together
still curve exactly like the simplex. Sampled sectional curvatures agree with $\tfrac14$ to
$10^{-18}$ over random natural parameters and $2$--planes (routine validated on four metrics of known
curvature); the family is not totally geodesic ($\log^2\Mah\notin V$). \forcedtag/\computedtag
\end{corollary}

\begin{theorem}[Geodesic threshold]\label{thm:threshold}
A family $V\ni\log\Mah$ is totally geodesic in $S^6(2)$ iff $V$ is a unital subalgebra of $\R^7$.
Any such $V$ contains $\R[\log\Mah]$, which is six--dimensional and equals the algebra of functions
constant on the golden pair. Hence no totally geodesic family with $k\le4$ statistics exists; the
first appears at exactly $k=5$, uniquely $V=\R[\log\Mah]$---the \emph{effective--catalog simplex}, a
great $S^5$ realizing the golden merge as forced geometry---and $k=6$ is the full simplex.
\forcedtag
\end{theorem}

\section{The $k$--window calculus and the $k=3$ census}\label{sec:kwindow}
The master identity lifts to every $k$ by the same three steps, and the lift carries a complete
classification of a first $k=3$ landscape.

\begin{theorem}[Generalized window identity]\label{thm:kmaster}
Let $T_1,\dots,T_k$ be statistics with $\dim V=k+1$, $V=\mathrm{span}(\mathbf 1,T)$, and for
$|s|=k+2$ put $\ell_s(f)=\det[(\mathbf 1,T_1,\dots,T_k,f)|_s]$. With
$C(u,v)=Z\langle u,v\rangle_w-\langle u,\mathbf 1\rangle_w\langle\mathbf 1,v\rangle_w$,
Sylvester's identity with pivot $Z$ and Cauchy--Binet give, exactly as at $k=2$,
\[
\det\bigl[C(u_i,v_j)\bigr]_{(k+1)\times(k+1)}=Z^{k}\!\!\sum_{|s|=k+2}\!\!w_s\,\ell_s(f)\ell_s(g),
\qquad
D_k\,\langle Q f,\,Q g\rangle_p=Z^{\,k-2}\!\!\sum_{|s|=k+2}\!\!w_s\,\ell_s(f)\ell_s(g),
\]
over rows $u=(f,T)$, columns $v=(g,T)$, with $D_k=\det[C(T_a,T_b)]=Z^{k-1}\det\mathrm{Gram}_w(\mathbf 1,T)$
and $Q$ the $L^2(p)$--projection off $V$. \forcedtag\ (\texttt{t4\_kwindows.py}: fully symbolic at
$k=3$, $m=5$; exact rational instances $m=6,7$; both Sylvester scalings.)
\end{theorem}

\begin{proposition}[Windowed Gauss obstructions and the flat criterion]\label{prop:kgauss}
The curvature--$\tfrac14$ obstructions
$G^{abcd}=\langle Q(T_aT_c),Q(T_bT_d)\rangle_p-\langle Q(T_aT_d),Q(T_bT_c)\rangle_p$ satisfy
\[
D_k\,G^{abcd}=Z^{\,k-2}\!\!\sum_{|s|=k+2}\!\!w_s\bigl[\ell_s(T_aT_c)\ell_s(T_bT_d)-\ell_s(T_aT_d)\ell_s(T_bT_c)\bigr],
\]
the right side running over the $2\times2$ minors of the \emph{window matrix}
$L(s)=[\ell_s(T_aT_b)]$. If $\dim V=k+1$ then $h\in V\iff\ell_s(h)=0$ for every window (flat
lemma). Hence: \emph{if every window matrix has rank $\le1$, all Gauss obstructions vanish
identically and the family has constant sectional curvature $\tfrac14$.} \forcedtag
\end{proposition}

\begin{corollary}\label{cor:kcatalog}
The catalog three--fold $(\log\Mah,\mathbf 1_K,\mathbf 1_{\ph^4})$ has all $21$ window matrices of
rank $\le1$ identically over $\Q(\log2,\log3,\log5,\log\ph)$, and the four--fold
$(\log\Mah,\mathbf 1_K,\mathbf 1_{\ph^4},\mathbf 1_{\sqrt5})$ has all $7$: exact linear--algebra
reconfirmations of Corollary~\ref{cor:threefold} and of Theorem~\ref{thm:join} at $k=4$, with no
curvature routine involved. \forcedtag
\end{corollary}

\begin{theorem}[The $k=3$ double--indicator census]\label{thm:kcensus}
Consider triples $(\log\Mah,\mathbf 1_S,\mathbf 1_T)$ over pairs of nonempty proper subsets with
$\dim V=4$ ($7812$ of the $7875$ unordered pairs). Windowwise flatness reduces, by rank $\le1$ on
the symmetric window matrix and the flat lemma, to three exact memberships
$\{a\mathbf 1_S,\;a\mathbf 1_T,\;\mathbf 1_{S\cap T}\}\subset V$ over $\Q(\log2,\log3,\log5,\log\ph)$.
Exactly $312$ pairs pass, and they collapse onto exactly $\mathbf{26}$ distinct spans
\[
V=\langle\mathbf 1,\;\log\Mah,\;\mathbf 1_{S'},\;\mathbf 1_{T'}\rangle,\qquad
S',T'\ \text{disjoint subsets each inside a single cost level}
\]
($\binom82-2$: eight within--level subsets, the two golden overlaps being redundant presentations
of the all--golden class, where $\mathbf 1_\ph+\mathbf 1_\tau=\mathbf 1_{\ph\tau}$ collapses). Each
of the $26$ families is constant--curvature $\tfrac14$ by Proposition~\ref{prop:kgauss}, and the
eight $k=2$ surfaces embed as the pairs with one within--level generator. \forcedtag\ (Two lanes:
\texttt{t4\_kwindows.py} decides membership by symbolic rank over the coefficient field;
\texttt{t4b\_census\_fast.py} re--derives the block--collapsed linear system and decides by exact
$\Q^5$--vector annihilator tests with quadratic cross--ratio consistency; sampled cross--lane
agreement on $24$ pairs $\times$ $3$ targets.)
\end{theorem}

\begin{remark}\label{rem:kobstruction}
What does \emph{not} lift verbatim is the $k=2$ chain ``windows $\Rightarrow$ trichotomy
$\Rightarrow$ plane lemma'': the multiset analog of the plane lemma is false at $k\ge3$ (a doubled
support point defeats the general--position projection). Both halves of the $k\ge3$ program
nonetheless close, by different routes: necessity in \S\ref{sec:knec} (squarefree collisions and
twin symmetry), and the landscape in \S\ref{sec:landscape2}, where the moment trichotomy makes
the doubled point harmless---a $2$--circuit never obstructs---and the plane lemma is bypassed
rather than repaired (Remark~\ref{rem:planebypass}). The census above is the double--indicator
stratum of the full landscape (Theorem~\ref{thm:paclass}).
\end{remark}

\begin{openproblem}\label{op:kclass}
\emph{Resolved in this edition, in both parts.} Necessity: Theorem~\ref{thm:knec}---constant
$\tfrac14$ is exactly windowwise rank $\le1$. Landscape: \S\ref{sec:landscape2}---constant
$\tfrac14$ is exactly the partitioned--affine condition $V=V(\pi)$, with catalog counts
$8/56/95/31/1$ for $k=2,\dots,6$; the conjecture of the previous edition, that within--level
indicator joins exhaust, is \emph{refuted} by the $30$ split--affine classes
$\langle1,a,\mathbf 1_B,a\mathbf 1_B\rangle$ (Corollary~\ref{cor:landres}).
\end{openproblem}

\section{Catalog invariance: the census theorem and the branch--rank calculus}\label{sec:census}
The eight--surface census is not an artifact of the seed selection: its shape is a theorem about
cost coincidences, and the coincidences are localized exactly.

\begin{theorem}[Generic census]\label{thm:generic}
If the costs $a_1,\dots,a_m$ are $\Q$--linearly independent, the collision lattice
$\{d\in\Z^m:\sum d_ia_i=0,\ \sum d_i=0\}$ is zero; by the master identity the vanishing $q\equiv0$
then forces every monomial coefficient of $P$ individually, hence every window $q_4(s)=0$, hence
(trichotomy and plane lemma, both catalog--free) exactly the $m$ single--seed surfaces. No branch
analysis occurs. \forcedtag
\end{theorem}

\begin{theorem}[Census count and catalog invariance]\label{thm:countform}
For any catalog, every within--level indicator $\mathbf 1_S$ ($S$ inside one cost level) is a
nondegenerate constant--$\tfrac14$ surface: in the factorization $q_4=-\prod\Delta$ each window
carries a vanishing factor (a coincident pair if $|s\cap S|\ge2$, three collinear points on $X=0$
otherwise). When the necessity analysis closes, the census is exactly
\[
\#\{\text{constant--}\tfrac14\text{ surfaces}\}\;=\;\sum_{\text{levels }\ell}\bigl(2^{m_\ell}-1\bigr),
\]
so the ``eighth surface'' is precisely the golden coincidence's $2^2-1=3$. Verified end to end for
four catalogs: the full catalog (collision lattice of rank $2$: golden swap $\Delta_1$, Salem square
$\Delta_2$; census $8$), drop--$K$ (rank $1$, golden; $7$), drop--$\tau$ (rank $1$, Salem; $6$), and
add--$\sqrt7$ (rank $2$ embedded; $9$). \forcedtag\ (\texttt{t5\_catalog\_census.py}, $72/72$.)
\end{theorem}

\begin{theorem}[Branch--rank calculus]\label{thm:branchrank}
Necessity closes for an arbitrary catalog by a rank criterion. The realizable primitive relation
directions $D$ are the primitive lattice vectors $d$ with $k$ and $k+d$ both monomials of $P$ (for
the full catalog, $D=\{\Delta_1,\Delta_2,\Delta_1\pm\Delta_2,2\Delta_1\pm\Delta_2\}$); the
branches are the flats of the arrangement they span. Per branch, group the monomials of
$P=Z^2\sum_sq_4(s)w_s$ by (cost class, $X$--class mod the active relations): $q\equiv0$ forces each
grouped $\Z$--combination of window symbols to vanish, and
\[
\operatorname{rank}_{\Q}(M_{\mathrm{branch}})=\#\{\text{windows}\}\;\Longrightarrow\;
\text{every }q_4(s)\text{ forced to }0 .
\]
A coincidence direction $\pm(e_i-e_j)$ with $a_i=a_j$ merges the two seeds as plane points: windows
containing both are automatically zero (coincident pair), the remaining symbols identify pairwise,
and the rank test runs against the $\binom{m-1}{4}$ surviving classes. This upgrades the $k=2$
singleton--survival lemma of \S\ref{sec:classdone} to a complete criterion (a unit heavy row per
window is the special case), mechanically reproduces the eight--branch tree of the full catalog, and
closes every branch of all four catalogs, with a one--sided mod--$p$ rank certificate and an exact
$\Q$ fallback. \forcedtag
\end{theorem}

\section{$k$--necessity: constant $\tfrac14$ is exactly windowwise flatness}\label{sec:knec}

Remark~\ref{rem:kobstruction} left necessity at $k\ge3$ open: the multiset plane lemma fails, so
the $k=2$ route does not lift. It turns out necessity never needed that route. The observation is
that the exponential identity behind the Gauss obstructions lives at the \emph{window} level,
where every monomial is squarefree---and the catalog's collision lattice has only one squarefree
direction.

\begin{theorem}[Window collisions are golden]\label{thm:goldencollide}
Let $T=(\log\Mah,T_2,\dots,T_k)$ be any statistics on the catalog, $2\le k\le5$. Two distinct
$(k{+}2)$--windows $s\ne s'$ have equal frequency vectors $\sum_{i\in s}T_i=\sum_{i\in s'}T_i$ iff
$s'=s\,\triangle\,\{\ph,\tau\}$ and \emph{every} statistic ties on the golden pair,
$T_{j,\ph}=T_{j,\tau}$ for all $j$. \forcedtag
\end{theorem}

\begin{proof}
The difference $\delta=\mathbf 1_s-\mathbf 1_{s'}$ has entries in $\{-1,0,1\}$ and zero sum, and
the cost coordinate forces $\delta$ into the zero--sum collision lattice $\langle d_1,d_2\rangle$
of Proposition~\ref{prop:pencil} ($d_1$ the golden swap, $d_2$ the Salem square). The general element
$jd_1+kd_2=(0,0,-k,\,j,\,-j,\,-k,\,2k)$ carries the entry $2k$ at $K$, so a cube point forces
$k=0$ and $|j|\le1$: $\delta=\pm d_1$, i.e.\ $s'=s\,\triangle\,\{\ph,\tau\}$. The swap then
preserves the remaining frequency coordinates iff every statistic ties on the pair.
\end{proof}

\begin{theorem}[$k$--necessity; the equivalence]\label{thm:knec}
On the catalog, for $2\le k\le 5$ and statistics $T=(\log\Mah,T_2,\dots,T_k)$ with $\dim V=k+1$:
the family has constant sectional curvature $\tfrac14$ iff every window matrix
$L(s)=[\ell_s(T_aT_b)]$ has rank $\le1$. \forcedtag\ (\texttt{t7\_knecessity.py}, $28/28$.)
\end{theorem}

\begin{proof}
Sufficiency is Proposition~\ref{prop:kgauss}. For necessity, constant $\tfrac14$ means the Gauss
obstructions $G^{abcd}$ of Proposition~\ref{prop:kgauss} vanish identically: they form an algebraic
curvature tensor, so vanishing of all sectional values forces the full tensor by polarization.
Since $D_k>0$ and $Z>0$ on the family, Proposition~\ref{prop:kgauss} turns $G^{abcd}\equiv0$ into the
exponential identity
\[
\sum_{|s|=k+2} w_s(\theta)\,M^{abcd}(s)\;\equiv\;0,\qquad
M^{abcd}(s)=\ell_s(T_aT_c)\,\ell_s(T_bT_d)-\ell_s(T_aT_d)\,\ell_s(T_bT_c),
\]
whose coefficients range, over all quadruples, over exactly the $2\times2$ minors of the symmetric
matrix $L(s)$ (rows $\{a,b\}$, columns $\{c,d\}$). If some statistic separates the golden pair,
Theorem~\ref{thm:goldencollide} gives pairwise distinct frequencies, so every $M^{abcd}(s)=0$
individually. If every statistic ties, the pair are \emph{twins}: for a swap pair the bordered
matrices $(\mathbf 1,T,T_aT_b)|_s$ and $(\mathbf 1,T,T_aT_b)|_{s'}$ agree row for row (the twins
occupy the same sorted slot, the indices being adjacent), so
$\ell_{s\triangle\{\ph,\tau\}}=\ell_s$ exactly and the grouped coefficient is $2M^{abcd}(s)$,
forcing each; and a window containing both twins carries two identical rows, so $L(s)=0$ outright.
In every case all $2\times2$ minors of every $L(s)$ vanish, i.e.\ rank $L(s)\le1$.
\end{proof}

\begin{corollary}[Open Problem~\ref{op:kclass}(i) closed; the census upgraded]\label{cor:knecclosed}
Constant $\tfrac14$ and windowwise flatness coincide unconditionally, so
Theorem~\ref{thm:kcensus}'s $26$ classes are exactly the constant--$\tfrac14$ double--indicator
triples---the complete such landscape---and Corollary~\ref{cor:kcatalog}'s three-- and four--folds
need no curvature routine. The analytic classification problem is reduced, at every $k$, to
finite linear algebra. \forcedtag
\end{corollary}

\begin{remark}[Squarefreeness; the branch tree retires]\label{rem:squarefree}
Windows are squarefree monomials; the Salem square is a multiplicity--two relation and cannot act
at window level---the same mechanism that made the heavy cofactor exactly $1$ in
Corollary~\ref{cor:windowlaw}(b). At $k=2$ this collapses the eight--branch analysis to the golden
dichotomy above: the branches lived on the \emph{expanded} monomials of $P$, where the prefactor
$Z^2$ reintroduces multiplicities; on the family one divides $Z^2$ out first. The branch machinery
(\S\ref{sec:classdone}, Theorem~\ref{thm:branchrank}) is retained as an independent verification lane
and as the general--catalog tool.
\end{remark}

\begin{theorem}[General catalogs]\label{thm:tame}
Call a catalog \emph{transposition--tame} if its zero--sum cost kernel meets $\{-1,0,1\}^m$ only
in signed sums of disjoint equal--cost transpositions. For such catalogs the equivalence of
Theorem~\ref{thm:knec} holds at every $k$, by the same twin argument classwise. All four studied
catalogs qualify: full, drop--$K$, and add--$\sqrt7$ (cube points $\pm d_1$ only), and
drop--$\tau$ (no cube points at all---no window collisions whatsoever). A catalog with a genuine
$\{-1,0,1\}$ trade (e.g.\ costs with $a_1+a_4=a_2+a_3$, all distinct) falls instead to the
window--level branch--rank test---group windows by frequency class and require rank
$=\#$windows---the $k$--window analog of Theorem~\ref{thm:branchrank}, strictly smaller than its
monomial--level version. \forcedtag
\end{theorem}

What remained of Open Problem~\ref{op:kclass}---its part (ii), the landscape beyond double
indicators---is closed in \S\ref{sec:landscape2}: the answer is the partitioned--affine
classification, and the conjecture recorded below it fails.

\section{The compositum couples: forced interaction between fields}\label{sec:compositum}

Open Problem (4) asked whether the compositum's own arithmetic forces a nonzero cross--structure
between two fields, or whether they remain geometrically independent even there. It resolves on
the first horn, with an exact witness.

For catalog seeds $p,q$ let $p\otimes q$ be the \emph{tensor polynomial}, whose roots are all
conjugate products $\alpha_i\beta_j$---no root is chosen, so the construction is gauge--free---and
let $c(p,q)=\log\Mah(p\otimes q)$ be the compositum cost. Every conjugate log--magnitude of the
catalog is a $\Q$--vector over $(\log2,\log3,\log5,\log\ph)$: for the quartic,
$|K_{\mathrm{out}}|^2=\sqrt5\,\ph^{2}$ and $|K_{\mathrm{in}}|^2=\sqrt5\,\ph^{-2}$. Hence each
$c(p,q)$ is an exact $\Q$--vector; products on the unit circle (zero vector, e.g.\ $\ph\psi=-1$)
are excluded symbolically, and every remaining sign is certified by rigorous interval arithmetic.

\begin{theorem}[The compositum couples]\label{thm:coupled}
The compositum cost is not additively separable: the interaction contrasts
\[
c(\sqrt2,\sqrt3)+c(\ph,\ph)-c(\sqrt2,\ph)-c(\ph,\sqrt3)=\log 2,
\qquad
c(\ph,\ph)+c(\ph^4,\ph^4)-2\,c(\ph,\ph^4)=-6\log\ph,
\]
both exactly, as nonzero coefficient vectors. Consequently the compositum's Mahler--Gibbs family
$p\propto\Mah(p\otimes q)^{-\beta}$ factors through the marginal families for no $\beta\ne0$; its
Fisher geometry carries nonzero cross--structure,
$\mathrm{Cov}_\beta(\mathrm{cost}_1,\mathrm{cost}_2)\ne0$ at the marked temperatures
\computedtag; and the coupling left \declaredtag\ by the product construction is, in the
compositum reading, \emph{forced}. \forcedtag\ (\texttt{t8\_compositum.py}, $13/13$.)
\end{theorem}

\begin{proposition}[Structure of the coupling]\label{prop:coupstruct}
(i) The rational sector $\{\sqrt2,\sqrt3,\sqrt5\}$ is exactly separable:
$c(\sqrt D,\sqrt{D'})=2\log D+2\log D'$, and all its interaction contrasts vanish. (ii) The
golden tie persists: $c(\ph,\cdot)=c(\tau,\cdot)$ (the conjugate magnitude multisets agree), so
the compositum inherits the golden merge. (iii) The double--centered interaction tensor $\Delta$
is nonzero of exact rank $5$ over $\Q(\log2,\log3,\log5,\log\ph)$: five independent coupling
modes. The coupling is a golden/quartic phenomenon; the fields do not decouple anywhere off the
rational sector. \forcedtag
\end{proposition}

\begin{proposition}[Charge tensor law, instance--verified]\label{prop:chargelcm}
With seed charge groups ($\Z/2$ at $\sqrt2,\sqrt3,\sqrt5,\ph,\tau$; $\Z/1$ at $\ph^4$; $\Z/4$ at
$K$), every tensor conjugate is real or purely imaginary, and on all $28$ unordered pairs the
tensor polynomial's charge group is exactly $\Z/\mathrm{lcm}(n_1,n_2)$---e.g.\
$\sqrt2\otimes K\to\Z/4$, $\ph^4\otimes\ph^4\to\Z/1$---consistent with the tensor charge law of
the charge--measure companion. \forcedtag\ (finite instance check)
\end{proposition}

\begin{remark}[Relation to the product geometry]\label{rem:coupwarp}
Proposition~\ref{prop:product} stands: the \emph{declared} product of two catalogs is a Riemannian
product with zero cross--curvature. The theorem says the compositum does not instantiate that
construction: its forced costs are not the product costs, and the discrepancy is the exact tensor
$\Delta$. What the product section left as a declared warp is here a forced nonzero datum. The two
questions this leaves---the coupled family's geometry, and whether $\Delta$ is
charge--determined---are settled in \S\ref{sec:coupledcurv}.
\end{remark}

\section{The landscape completed: the partitioned--affine classification}\label{sec:landscape2}

With Theorem~\ref{thm:knec}, constant $\tfrac14$ \emph{is} windowwise flatness; what Open
Problem~\ref{op:kclass} still asked for was the classification of the windowwise--flat spans
$V\ni 1,a$ with $\dim V=k+1$. This section completes it, at every $k$, and the conjecture
recorded there turns out to be \emph{false}. The right chart is the point configuration: write
$P_i=(T_1,\dots,T_k)(i)\in\R^k$ for the statistic values at seed $i$ (with $T_1=a$). For a
window $s$ whose points affinely span $\R^k$, the functional $\ell_s$ is exactly the unique
affine dependency $\lambda$ of the $k+2$ points---supported on the window's unique circuit---and
\[
L(s)\;=\;\bigl[\ell_s(T_aT_b)\bigr]_{a,b}\;=\;\sum_{i\in s}\lambda_i\,P_iP_i^{\top},
\]
whose rank is invariant under affine changes of the statistic basis (the cross terms die on
$\sum\lambda_i=0$, $\sum\lambda_iP_i=0$). Non--spanning windows carry $L(s)=0$ and no condition,
so the window criterion is a statement about circuits.

\begin{theorem}[Moment trichotomy]\label{thm:momtri}
Let $s$ be a spanning window with unique circuit $C$ and dependency $\lambda$. If $|C|=2$ (a
coincident pair), $L(s)=0$. If $|C|=3$ (a collinear triple with parameters $t_1,t_2,t_3$),
$L(s)$ has rank exactly $1$, with moment
$\sum_i\lambda_it_i^2=-(t_1-t_2)(t_2-t_3)(t_3-t_1)\ne0$. If $|C|\ge4$,
$\operatorname{rank}L(s)\ge2$. Consequently a $k$--family is constant--$\tfrac14$ if and only if
every circuit of its configuration has size $\le3$. \forcedtag
\end{theorem}

\begin{proof}
Only $|C|\ge4$ needs an argument. A circuit's points are in general position in their
$(d{=}|C|{-}2)$--flat; affinely normalize $d+1$ of them to $0,e_1,\dots,e_d$, so the last is
$(x_1,\dots,x_d)$ and $\lambda=(\sum_jx_j-1,\,-x_1,\dots,-x_d,\,1)$, all entries nonzero on a
circuit. Then $Q=\sum\lambda_iP_iP_i^{\top}$ has $Q_{jj}=x_j^2-x_j$ and $Q_{jl}=x_jx_l$. For
$d\ge3$ the mixed $2\times2$ minor on rows $\{j,l\}$, columns $\{l,r\}$ equals $x_jx_lx_r$; for
$d=2$, $\det Q=xy(1-x-y)$. Both are products of circuit weights, hence nonzero.
(\texttt{t9}, blocks 1--2: symbolic at $d=2,3,4$, plus exact instances.)
\end{proof}

\begin{theorem}[Structure: flat $=$ partitioned--affine]\label{thm:pastruct}
Suppose the configuration is windowwise flat. Then its maximal collinear location--sets with
$\ge3$ locations (\emph{rich lines}) are pairwise disjoint and pairwise skew; every collinear
triple of locations lies inside a single rich line; and, writing $\pi$ for the partition of the
seeds into rich--line blocks and free locations (coincident clusters and singletons),
\[
V\;=\;V(\pi)\;:=\;\{\,f:\ f\ \text{is affine in }a\text{ on every block of }\pi\,\}.
\]
Conversely, every $V(\pi)$ with $\dim V(\pi)=k+1$ is windowwise flat. \forcedtag
\end{theorem}

\begin{proof}
Two rich lines sharing a location or lying in a common plane yield four coplanar locations, two
per line, with no collinear triple---a $4$--circuit---contradicting Theorem~\ref{thm:momtri}; so
rich lines are disjoint and skew. A collinear triple across structures would make its own line
rich and meeting another. On each rich line the cost is affine and nonconstant (three locations
in one level would need level multiplicity $\ge3$), hence an injective parameter, and every
statistic is affine along the line; so $V\subseteq V(\pi)$. If $\dim V<\dim V(\pi)=2r+f$, the
affine hull is deficient and a transversal selection---two locations per line, one per free
location---is affinely dependent; its circuit has no coincidence, at most two points per line,
and no transversal triple, hence size $\ge4$: again a contradiction. So $\dim V=2r+f$ and
$V=V(\pi)$. For sufficiency, in the block chart the functionals $\mathbf 1_{B_j}$,
$a\mathbf 1_{B_j}$, $\mathbf 1_{C}$ annihilate any dependency using $\le2$ points per line block
and $\le1$ per free location (two locations of a line block carry distinct costs), so every
circuit is a coincident pair or an in--block collinear triple, and Theorem~\ref{thm:momtri}
gives rank $\le1$ on every spanning window.
\end{proof}

A block with at most two distinct cost values is the \emph{same} function space as its
within--level refinement---affine in $a$ through two values is no constraint---so the canonical
presentation has two block types: \emph{clusters}, within--level subsets (on this catalog
singletons or the golden pair $\{\ph,1/\ph\}$), and \emph{line blocks} carrying at least three
distinct cost values. Then $\dim V(\pi)=c+2b_2$ with $c$ clusters and $b_2$ line blocks.

\begin{theorem}[Classification and counts]\label{thm:paclass}
The constant--$\tfrac14$ $k$--families of the catalog are exactly the canonical $V(\pi)$ with
$c+2b_2=k+1$, and the counts are
\[
k=2{:}\ 8,\qquad k=3{:}\ 56,\qquad k=4{:}\ 95,\qquad k=5{:}\ 31,\qquad k=6{:}\ 1 .
\]
At $k=2$ this recovers the eight surfaces of Theorem~\ref{thm:classification} in one line. At
$k=3$ the $26$ double--indicator classes of Theorem~\ref{thm:kcensus} are joined by $30$
\emph{split--affine} classes $\langle 1,a,\mathbf 1_B,a\mathbf 1_B\rangle$---the two--block
partitions with at least three cost values on both sides; the five $(3,4)$--splits whose
$3$--side is $\{\ph,1/\ph,x\}$ coincide with $\langle1,a,\mathbf 1_{\{\ph,1/\ph\}},\mathbf
1_x\rangle$ and are not new. At $k=5$ the golden merge $\R[\log M]$---the totally geodesic
family of Theorem~\ref{thm:threshold}---appears among the $31$; at $k=6$ only the full simplex
remains. \forcedtag\ (\texttt{t9\_landscape.py}: exhaustive exact window verification of all
$8+56+95+31$ families over $\Q(\log2,\log3,\log5,\log\ph)$; the five overlap equalities by
symbolic rank; pairwise distinctness of the $56$ certified at the true logarithms by $1540$
interval rank--$5$ witnesses; twenty exact random controls violate flatness.)
\end{theorem}

\begin{corollary}[Open Problem~\ref{op:kclass} resolved; the conjecture refuted]\label{cor:landres}
Within--level indicator joins do \emph{not} exhaust the landscape: a split--affine class
contains no indicators beyond $\mathbf 1_B$ and $\mathbf 1_{B^c}$ (an affine $0/1$ pattern
across three or more distinct cost values is constant), so it is not an indicator join, yet it
is constant--$\tfrac14$. Theorem~\ref{thm:kcensus} stands exactly as the double--indicator
stratum of the $56$. \forcedtag
\end{corollary}

\begin{remark}[The plane lemma, bypassed]\label{rem:planebypass}
The moment trichotomy replaces the $k=2$ chain wholesale. The multiset failure of
Remark~\ref{rem:kobstruction} is not repaired but made irrelevant: a doubled support point is a
$2$--circuit and never obstructs. At $k=2$ the structure theorem \emph{implies} the plane lemma
(the plane holds no two skew lines, so $r\le1$ and $f\le1$), which is how the eight surfaces
drop out.
\end{remark}

\begin{remark}[Scope: level multiplicities]\label{rem:fatlevels}
The proofs use only that every cost level of the catalog holds at most two seeds; the theorem
applies verbatim to the drop--$K$, drop--$1/\ph$ and add--$\sqrt7$ catalogs. A catalog with a
level of multiplicity $\ge3$ admits genuinely new \emph{vertical} line blocks (a line inside one
level, transverse to the cost coordinate), classified by the same proof with the line parameter
replacing $a$; the counts above are for the studied catalogs.
\end{remark}

\section{The coupled family curves}\label{sec:coupledcurv}

Remark~\ref{rem:coupwarp} left two residuals: the geometry of the coupled family
$(\mathrm{cost}_1,\mathrm{cost}_2,c)$ on the $49$--cell tensor catalog, and whether $\Delta$ is
determined by the charge data. Both close, the first with the machinery of
\S\ref{sec:landscape2}.

\begin{theorem}[The coupled family is not constant--$\tfrac14$]\label{thm:coupcurv}
On the tensor catalog the triple $T=(\mathrm{cost}_1,\mathrm{cost}_2,c)$ spans, $\dim V=4$, and
its configuration contains a genuine $5$--circuit---all five dependency weights nonzero,
certified at the true logarithms by interval arithmetic---so by Theorem~\ref{thm:pastruct} it is
\emph{not} of partitioned--affine type, and the corresponding window matrix has rank $\ge2$: the
family is not windowwise flat. Directly, the Gauss obstructions at the uniform point are nonzero
(rigorous enclosures), with sectional curvatures
\[
\sec(\mathrm{cost}_1,\mathrm{cost}_2)\;=\;\tfrac14-0.248531\ldots\;\approx\;0.0015,
\qquad
\sec(\mathrm{cost}_i,c)\;=\;\tfrac14-0.083802\ldots\;\approx\;0.1662 .
\]
The forced coupling of Theorem~\ref{thm:coupled} is curvature--visible; the near--vanishing of
the first plane is the ghost of the product geometry (Proposition~\ref{prop:product}).
\forcedtag/\computedtag\ (\texttt{t10\_coupled.py}, $9/9$; enclosures at $60$ digits.)
\end{theorem}

\begin{proposition}[Invisible to covariance, visible to curvature]\label{prop:covinvis}
At the uniform point $\Cov(\mathrm{cost}_1,\mathrm{cost}_2)=0$ exactly---the marginals are
independent there---yet the sectional curvatures above already deviate from $\tfrac14$. The
coupling enters through the compositum cost $c$, below the reach of the cross--covariance at
that point: the covariance witness of Theorem~\ref{thm:coupled} lives at $\beta\ne0$, the
curvature witness already at $\beta=0$. \forcedtag
\end{proposition}

\begin{proposition}[$\Delta$ is not determined by the charge data]\label{prop:deltacharge}
The pairs $(\sqrt2,\sqrt3)$ and $(\sqrt2,\sqrt5)$ carry the same charge pair $(\Z/2,\Z/2)$, but
\[
\Delta(\sqrt2,\sqrt3)=\tfrac{5}{49}\log2-\tfrac{5}{49}\log3+\tfrac{20}{49}\log\ph
\;\ne\;
\Delta(\sqrt2,\sqrt5)=\tfrac{5}{49}\log2+\tfrac{2}{49}\log3-\tfrac17\log5+\tfrac{20}{49}\log\ph,
\]
exactly. The interaction tensor is a forced datum beyond the charge grading. \forcedtag
\end{proposition}

\begin{remark}[The ledger closes]\label{rem:closed}
With \S\ref{sec:landscape2} and this section, every enumerated open problem of the corpus is
resolved. The temperature selection remains \declaredtag{} by the dichotomy of
Theorem~\ref{thm:dichotomy}---its honest final form---and the standing conditioning is, as
throughout, the $\Q$--linear independence of the seed logarithms.
\end{remark}

\section{Interaction with the trace form}
\begin{theorem}[Embedding obstruction]\label{thm:embed}
The field's flat trace--form geometry admits no totally geodesic isometric embedding into the
surprisal sphere; any isometric surface realization is forced to satisfy $k_1k_2=-\tfrac14$
(Gauss: $0=\tfrac14+k_1k_2$), witnessed by the Clifford torus in $S^3(2)$ with principal curvatures
$\pm\tfrac12$. \forcedtag
\end{theorem}

\begin{theorem}[Fisher monotonicity]\label{thm:mono}
Coarse--graining the field's conjugate data into the catalog outcomes is a Markov morphism under
which $I_{\mathrm{surprisal}}\preceq I_{\mathrm{trace}}$ (data--processing \cite{zamir,chentsov}),
equality iff sufficient. The catalog surprisal geometry is therefore a lossy, one--way contraction
of the trace--form geometry. \forcedtag
\end{theorem}

Together with Theorem~\ref{thm:susp} the picture is sharp: the arithmetic (trace--form) geometry is
flat and signature--flexible; the surprisal geometry is curved ($\tfrac14$) and positive definite;
and the two--statistic field surface sits inside the latter as a constant--curvature--$\tfrac14$ ruled
surface---neither it nor the trace form is totally geodesic (Prop.~\ref{prop:notTG}), for the opposite
reasons of wrong embedding and wrong intrinsic curvature respectively.

\section{Epistemic ledger}
\begin{table}[h]\centering
\begin{tabular}{lc}
\toprule
claim & status\\ \midrule
$Z=17+4\sqrt5$; Mahler--Gibbs is an exponential family with statistic $\log\Mah$ & \forcedtag\\
surprisal $=\beta\log\Mah+\log Z$; Fisher $=\Var_\beta(\log\Mah)$; geometry $=S^{m-1}(2)$, $K=\tfrac14$ & \forcedtag\\
dual flatness; $\mathrm{KL}$ $=$ Bregman divergence of $\log Z$; $A^\ast=-H$ & \forcedtag\\
$C(\beta)=\beta^2 I(\beta)$ (Fisher information is a heat capacity) & \forcedtag\\
$d_{\mathrm{FR}}=2\arccos$ BC; vertices mutually $\pi$ apart & \forcedtag\\
keystone sequence: shifted geometric law; $I=(\log\ph)^2 r/(1-r)^2$ & \forcedtag\\
embedding obstruction $k_1k_2=-\tfrac14$; Fisher monotonicity (lossy contraction) & \forcedtag\\
$(\log\Mah,\mathrm{deg})$ Fisher matrix SPD; peak--fluctuation $\beta^\ast$; information length & \computedtag\\
no $(\log\Mah,X)$ surface is totally geodesic ($\log\Mah$ spans a $6$--dim algebra) & \forcedtag\\
$(\log\Mah,\mathrm{deg})$ surface: constant curvature $\tfrac14$, ruled (indicator--suspension) & \forcedtag\\
among forced invariants only degree gives constant $\tfrac14$; $\#$outside/trace/house vary & \computedtag\\
product of independent catalogs: block--diagonal Fisher, zero cross--curvature & \forcedtag\\
charge parity $=\mathbf 1_{\{\ph^4\}}$: a second constant--$\tfrac14$ surface (dual to degree) & \forcedtag\\
exactly eight constant--$\tfrac14$ surfaces; non--indicator exclusion by branch calculus & \forcedtag\\
master window identity $P=Z^2\sum_sq_4(s)w_s$ (Sylvester $+$ Cauchy--Binet); heavy cofactor $1$ & \forcedtag\\
iterated join: multi--indicator families are round $\tfrac14$--spheres; catalog $3$--fold & \forcedtag\\
geodesic threshold $k=5$, uniquely $\R[\log\Mah]$ (the golden--merged simplex) & \forcedtag\\
total information length $L_{\mathrm{tot}}\approx5.6462$; endpoints golden merge $\leftrightarrow\ \ph^4$ & \computedtag\\
coupling of two fields is a declared warp (\v{C}encov), like $\beta$ and $c$ & \declaredtag\\
Gibbs form $p\propto\Mah^{-\beta}$ forced by product functoriality $+$ measurable Cauchy & \forcedtag\\
selection of a single operating temperature among the marked points & \declaredtag\\
generalized window identity, every $k$ (Sylvester $+$ Cauchy--Binet over $(k{+}2)$--subsets) & \forcedtag\\
window criterion: rank--$\le1$ windows $\Rightarrow$ constant $\tfrac14$; catalog $3$-- and $4$--folds exact & \forcedtag\\
$k=3$ double--indicator census: exactly $26$ windowwise--flat classes (disjoint within--level pairs) & \forcedtag\\
generic census: independent costs $\Rightarrow$ coefficientwise vanishing $\Rightarrow$ $m$ surfaces & \forcedtag\\
census count $\sum_\ell(2^{m_\ell}{-}1)$; catalogs full/drop--$K$/drop--$\tau$/add--$\sqrt7$: $8/7/6/9$ & \forcedtag\\
branch--rank calculus (grouped window symbols, rank $=\#$windows per branch) & \forcedtag\\
selection dichotomy: no metric--canonical point; midpoint/max--Fisher/max--entropy/peak--$C$ distinct & \forcedtag/\computedtag\\
$k$--necessity, every $k$: constant $\tfrac14$ $\Leftrightarrow$ windowwise rank $\le1$ (squarefree collisions $+$ twin symmetry) & \forcedtag\\
window collisions are golden--only: the Salem square cannot act at window level & \forcedtag\\
compositum forces coupling: contrast exactly $\log2$; rational sector separable; golden tie persists; $\mathrm{rank}\,\Delta=5$ & \forcedtag\\
charge tensor law $n(p\otimes q)=\mathrm{lcm}$ on all $28$ compositum pairs; cross--Fisher $\mathrm{Cov}_\beta\ne0$ & \forcedtag/\computedtag\\
moment trichotomy: window circuits of size $2/3/{\ge}4$ give rank $0/1/{\ge}2$ & \forcedtag\\
landscape classified: constant--$\tfrac14$ $=$ partitioned--affine $V(\pi)$; counts $8/56/95/31/1$, $k=2..6$ & \forcedtag\\
indicator--join conjecture refuted: $30$ split--affine classes $\langle1,a,\mathbf 1_B,a\mathbf 1_B\rangle$ & \forcedtag\\
coupled family not constant--$\tfrac14$: Gauss obstruction $\ne0$ at uniform; $\Delta$ not charge--determined & \forcedtag/\computedtag\\
$Z$'s integer $17$; $I(1/\ph)=\ph^3$ & {\scshape [coincidence]}\\
\bottomrule
\end{tabular}
\end{table}

\section{Open problems}
(1) \emph{Resolved in this edition as a dichotomy} (Theorem~\ref{thm:dichotomy}): no invariance
principle selects an operating point canonically---the labeled--endpoint isometry group is trivial,
and the candidate principles (arc--length midpoint, maximal Fisher information, maximum entropy,
peak heat capacity) select provably distinct temperatures. The \declaredtag{} on the selection is
the answer's honest form, not a gap.
(2) \emph{Resolved in this edition} (\S\ref{sec:landscape2}): with necessity closed at every
$k$ (Theorem~\ref{thm:knec}), constant $\tfrac14$ \emph{is} windowwise flatness, and the
windowwise--flat spans are classified---$V=V(\pi)$, partitioned--affine, with catalog counts
$8/56/95/31/1$ for $k=2,\dots,6$. The indicator--join conjecture is refuted by the $30$
split--affine classes at $k=3$ (Corollary~\ref{cor:landres}).
(3) \emph{Resolved in this edition.} The conceptual proof of the window identity is
Theorem~\ref{thm:master}: $P=Z^2\sum_{|s|=4}q_4(s)\,w_s$ by Sylvester's determinant identity (pivot
$Z$) plus Cauchy--Binet, with the heavy cofactor--$1$ law, the full coefficient law and $462$--census,
the factorization $q_4=-\prod\Delta$, and the compound--ratio form of $q$ as corollaries
(Corollary~\ref{cor:windowlaw}). What was found by machine ($140$ exact divisions) is now a two--line
determinantal mechanism; the collision/branch machinery is retired from the critical path
(Remark~\ref{rem:master}).
(4) \emph{Resolved in this edition, on the first horn} (\S\ref{sec:compositum}): the
compositum's own tensor catalog forces the coupling---the cost $\log\Mah(p\otimes q)$ is provably
non--separable, with interaction contrast exactly $\log 2$ (and $-6\log\ph$ in the golden
sector), nonzero cross--Fisher covariance, and an interaction tensor of exact rank $5$; the
tensor charge law $\Z/\mathrm{lcm}$ holds on all $28$ pairs. The declared warp is promoted to a
forced nonzero datum in the compositum reading. Its residuals are closed in
\S\ref{sec:coupledcurv}: the coupled family is provably not constant--$\tfrac14$---the coupling
is curvature--visible, already at the uniform point where the cross--covariance vanishes
exactly---and $\Delta$ is not determined by the charge data.
(5) \emph{Resolved in this edition} (Theorem~\ref{thm:knec}): constant $\tfrac14$
$\Rightarrow$ windowwise rank $\le1$ at every $k$, by squarefree window collisions (the Salem
square cannot act on windows) plus twin symmetry at the golden pair.

\medskip
\noindent\emph{Outlook.} What remains is outward--facing rather than internal: catalogs with
cost levels of multiplicity $\ge3$, where vertical line blocks enter the partitioned--affine
classification (Remark~\ref{rem:fatlevels}); composita of two \emph{distinct} catalogs, where
the tensor construction and $\Delta$ generalize verbatim; the coupled family's curvature
landscape away from the uniform point; and, as throughout, the standing conditioning of the
collision analysis on the $\Q$--linear independence of the seed logarithms.

\appendix
\section{Reproducible verification}\label{app:verify}
All load--bearing identities are checked in \texttt{field\_surprisal\_v2.py} (\texttt{sympy} exact
arithmetic; \texttt{mpmath} for transcendental and numeric values), printing $12/12$ passed. The
checks: $Z=17+4\sqrt5$, the affinity, and $I=\Var(\log\Mah)$ (\S\S3--4); the Bregman--$\mathrm{KL}$
identity and $A^\ast=-H$ (\S5); $C=\beta^2 I$ and the peak--fluctuation temperature (\S6); the
vertex distance $\pi$ and the information length (\S7); the finite geometric closed form and the
geometric variance of the keystone sub--model (\S9); the two--statistic Fisher matrix and the
peak--fluctuation temperature (\S\S6,10); and the embedding obstruction $k_1k_2=-\tfrac14$ and the
strict Fisher contraction (\S11). Theorem~\ref{thm:susp}---the exact constant curvature $\tfrac14$ of
the $(\log\Mah,\mathrm{deg})$ surface---is proved in the text and verified symbolically in
\texttt{suspension\_theorem.py} ($9/9$): the Brioschi curvature of a general indicator family equals
$\tfrac14$ identically in the free statistic, remains $\tfrac14$ under arbitrary base measure (hence
constant everywhere), the suspension metric has $g_{\psi\psi}=1,\ g_{\theta\psi}=0$ and curvature
$1$, the catalog specialization gives $\tfrac14$, and generic controls give $4/21$ and $857/16928$.
The Tier--2 results---the multi--statistic landscape (only the single--outcome indicator is
constant $\tfrac14$; the $3$--set indicator and generic statistics vary) and the product geometry
(block--diagonal Fisher, zero cross--curvature), and the charge grading (charge parity
$=\mathbf 1_{\{\ph^4\}}$ is a second constant--$\tfrac14$ surface; the charge--group order varies)---are
verified in \texttt{field\_surprisal\_tier2.py} ($15/15$). The complete classification---no surface
totally geodesic ($d\in\{6,7\}$, $\log\Mah^2\notin V$), the ruled second fundamental form
($\mathrm{II}_{ac}=\mathrm{II}_{cc}=0$, $\mathrm{II}_{aa}\neq0$), the exactly eight constant--$\tfrac14$
surfaces among all $127$ indicators, and the $\approx5000$--sample exclusion of non--indicators---is
verified in \texttt{field\_surprisal\_classification.py} ($9/9$). Across the four harnesses ($12/12$,
$9/9$, $15/15$, $9/9$) a reviewer reruns every value; all numbered results are proved in full or
established as marked.

The closures of \S\S\ref{sec:classdone}--\ref{sec:join} are verified in six further harnesses,
with exact \texttt{sympy} arithmetic over $\Q(\log2,\log3,\log5,\log\ph)$ at every decision
boundary and floats display--only: \texttt{t1\_core.py} ($61/61$: cost--value structure, the ruled
system, the invariant--module dimensions $N(X)$, the $r=\dim V^2/V$ census, the M\"obius elimination
steps); \texttt{t1\_reduction.py} ($17/17$: $Z^2Dq=P$, $q=\operatorname{Tr}(\Gamma N)$, the
four--point rank--one law $q=\kappa q_4$, the rank--$2$ collision lattice); \texttt{t1\_engine.py}
($4/4$: the full $462$--monomial coefficient dictionary of $P$, support--$\le3$ vanishing, the window
identity at all $140$ heavy placements, and the $\nu$--group census); \texttt{t1\_branches.py}
($16/16$: window survival in every branch of the collision pencil, the merged six--point recursion
with its rank--$1$ lattice, and identical window vanishing on the eight families in symbolic
$(\alpha,\beta,\gamma)$); \texttt{t2\_temperature.py} ($12/12$: factorization rank tests, the
lattice--Cauchy pathology, the exact anchors $Z(\mp1)$, the identity $I'=-\kappa_3$, and the curve
invariants); \texttt{t3\_suspension.py} ($24/24$: the join metrics' constant curvature for $k=3,4$
exactly, curvature routines validated on sphere, hyperbolic plane, flat plane, and the simplex, the
catalog three--fold at $\tfrac14$ to $10^{-18}$, and the geodesic threshold). Session total: $134$
checks, all exit $0$.

Two harnesses added in this edition. \texttt{t1\_windowproof.py} ($44/44$, exit $0$) machine--proves
Theorem~\ref{thm:master} and Corollary~\ref{cor:windowlaw}: the Sylvester and Cauchy--Binet steps as
fully symbolic polynomial identities at $m=4,5$; the factorization
$q_4=-\prod_{t\in s}\Delta_{s\setminus t}$ symbolically for four generic points; the compound--ratio
form of $q$ and the face law $q=\kappa q_4$ at exact rational instances for $m=6,7$; and, on the
catalog ring $\Q[X_1,\dots,X_7,\log2,\log3,\log5,\log\ph]$, the dictionary--level equality of all
$462$ coefficients of $P$ with $Z^2\sum_sq_4(s)w_s$, string--exact against the artifact saved by
\texttt{t1\_engine.py}. \emph{Ordering note}: that final cross--check reads the artifact
\texttt{P\_coeffs.pkl} written by \texttt{t1\_engine.py}, so run \texttt{t1\_engine.py} first for the
full $44/44$; standalone, the artifact check auto--skips (reported ``not counted'') and the tally is
$43/43$. The one--line runner \texttt{run\_all.sh} executes all harnesses in a correct order. \texttt{t3c\_partC\_exact.py} ($16/16$, exit $0$) re--proves the Part--C
claims of \texttt{t3\_suspension.py} (checks $[016]$--$[024]$) by explicit certificates that avoid
symbolic--rank heuristics: the level--set annihilator $\prod_{r=1}^{6}(a-v_r\mathbf 1)=0$ identically
in $\Q[\log2,\log3,\log5,\log\ph]^7$, the explicit closure recursion for $a^k$ ($k\ge6$) with
elementary--symmetric coefficients, the Vandermonde lower bound (conditioned, as everywhere, only on
the $\Q$--linear independence of the four logarithms), the golden--pair span equality, the $k\le4$
exclusion and $k=5$ uniqueness, and the non--geodesy of the catalog three--fold.

Four harnesses close this edition's remaining fronts. \texttt{t4\_kwindows.py} proves the
generalized window identity symbolically at $k=3$, $m=5$, checks it at exact rational instances
$m=6,7$ with both Sylvester scalings, verifies the windowed Gauss--obstruction identity and the flat
lemma, certifies the catalog three-- and four--fold window matrices rank $\le1$ over the coefficient
field, and runs the double--indicator census by symbolic--rank membership (checks
$[001]$--$[013]$, exit--clean in--sandbox; the census block's symbolic ranks exceed the execution
ceiling there and are superseded by the second lane).
\texttt{t4b\_census\_fast.py} ($9/9$, exit $0$) is the census's second lane: an independent
block--system reduction deciding the same memberships by exact $\Q^5$--annihilator and quadratic
cross--ratio tests, completing all $7875$ pairs ($7812$ nondegenerate, $312$ flat, $26$ classes)
with sampled cross--lane agreement. \texttt{t5\_catalog\_census.py} ($72/72$, exit $0$) proves the
generic census theorem, the count formula, the collision lattices (ranks $2/1/1/2$ with the named
generators, $\Z$--saturated), the sufficiency factorization, and the branch--rank necessity in every
branch of all four catalogs, the full catalog reproducing exactly the eight--branch tree of
\texttt{t1\_branches.py}. \texttt{t6\_selection.py} ($18/18$, exit $0$) proves the exact anchors
$Z(+1)=\tfrac{91}{30}-\tfrac{\sqrt5}5$ and $Z(-1)=17+4\sqrt5$ in $\Q(\sqrt5)$, the identity
$I'=-\kappa_3$ symbolically, and the selection dichotomy's numeric margins at dps $30/50$
agreement below $10^{-12}$.

Two harnesses close the final fronts. \texttt{t7\_knecessity.py} ($28/28$, exit $0$) is
the machine lane of \S\ref{sec:knec}: the window--collision census of all four catalogs at every
window size (every collision class is a single golden--swap pair; drop--$\tau$ has none), the
cube fact for $\langle d_1,d_2\rangle$, twin symmetry of the window functionals at $k=2$ (all ten
swap pairs and all ten twin--containing windows, by $4\times4$ determinants) and $k=3$ (entrywise
row agreement plus determinant spot--checks), the minor/rank equivalence, the $k=3$ windowed
identity $D_kG^{abcd}=Z^{k-2}\sum_sw_s[\,2\times2\ \text{minor}\,]$ at an exact rational instance
($m=7$, all nine quadruples), and exact windowwise--rank recertification of the catalog three--
and four--folds with a failing control. \texttt{t8\_compositum.py} ($13/13$, exit $0$) verifies
\S\ref{sec:compositum}: the $7\times7$ tensor--cost matrix as exact $\Q$--vectors (the ten
on--circle root products excluded symbolically; every remaining sign certified by rigorous
interval arithmetic), the two exact interaction witnesses, rational--sector separability, the
persistent golden tie, $\mathrm{rank}\,\Delta=5$ over $\Q(\log2,\log3,\log5,\log\ph)$, nonzero
cross--Fisher covariance at $\beta\in\{1,-1,\sqrt5\}$, and the $\mathrm{lcm}$ charge law on all
$28$ pairs. A practical note: the exhaustive \emph{symbolic} census lane of
\texttt{t4\_kwindows.py} (block 3) exceeds a sandbox execution ceiling; the fast
$\Q^5$--annihilator lane \texttt{t4b\_census\_fast.py} is the census's practical lane, with
sampled two--lane agreement.

Two final harnesses. \texttt{t9\_landscape.py} ($21/21$, exit $0$): the moment trichotomy
symbolically at $d=2,3,4$ and on exact instances; the full canonical enumeration with every
window matrix certified rank $\le1$ over $\Q(\log2,\log3,\log5,\log\ph)$ ($26+30$ classes at
$k=3$, $50+45$ at $k=4$, $30+1$ at $k=5$, the eight at $k=2$); the five golden--overlap span
equalities by symbolic rank; pairwise distinctness of the $56$ at the true logarithms ($1540$
rank--$5$ witness minors, interval--certified); and twenty exact random controls violating
flatness. \texttt{t10\_coupled.py} ($9/9$, exit $0$): the coupled triple's spanning
($\dim V=4$); a genuine $5$--circuit and a rank--$\ge2$ window at the true logarithms; the three
nonzero Gauss obstructions at the uniform point with certified sectional enclosures; the exact
vanishing of the uniform cross--covariance; and the $\Delta$--charge witness.

\begin{thebibliography}{9}
\bibitem{amari} S.~Amari and H.~Nagaoka, \emph{Methods of Information Geometry}, AMS/OUP, 2000.
\bibitem{rao} C.~R.~Rao, ``Information and accuracy attainable in the estimation of statistical
parameters,'' \emph{Bull. Calcutta Math. Soc.} \textbf{37} (1945), 81--91.
\bibitem{chentsov} N.~N.~\v{C}encov, \emph{Statistical Decision Rules and Optimal Inference}, AMS,
1982.
\bibitem{zamir} R.~Zamir, ``A proof of the Fisher information inequality via a data processing
argument,'' \emph{IEEE Trans. Inform. Theory} \textbf{44} (1998), 1246--1250.
\bibitem{rissanen} J.~Rissanen, ``Modeling by shortest data description,'' \emph{Automatica}
\textbf{14} (1978), 465--471.
\bibitem{shima} H.~Shima, \emph{The Geometry of Hessian Structures}, World Scientific, 2007.
\bibitem{aczel} J.~Acz\'el, \emph{Lectures on Functional Equations and Their Applications},
Academic Press, 1966.
\bibitem{hj} R.~A.~Horn and C.~R.~Johnson, \emph{Matrix Analysis}, 2nd ed., Cambridge University
Press, 2013.
\bibitem{smyth} C.~J.~Smyth, ``On the product of the conjugates outside the unit circle of an
algebraic integer,'' \emph{Bull. London Math. Soc.} \textbf{3} (1971), 169--175.
\bibitem{cmc} Echo--Squirrel Research, ``The Charge--Measure Coupling on a Spectral Semiring:
Cyclicity, Composition, and a Parity-Graded Mahler Floor,'' 2026.
\bibitem{lambda} Echo--Squirrel Research, ``The Exchange Rate $\lambda=2c$,'' companion manuscript,
2026.
\end{thebibliography}
\end{document}
